% !TeX program = pdflatex
% 基于提供的 Medical Physics 2021 模板。使用 pdfLaTeX 编译两次。
% 内容来自 manuscript-20261114.docx；修改正文后重新编译。
\documentclass[12pt,twoside,letterpaper]{article}
\usepackage[T1]{fontenc}
\usepackage[utf8]{inputenc}
\usepackage{lmodern}
\usepackage{amsmath,amssymb}            % 数学公式
\usepackage{graphicx}                   % 图片
\usepackage{placeins}                   % 在必要边界清理待排图表
\usepackage{float}                      % 图表保留在对应正文位置
\usepackage{booktabs,array,tabularx}                   % 三线表
\usepackage[super,sort&compress,comma]{natbib} % 上标数字引用，连续编号合并为范围
\usepackage{setspace}                   % 行距
\usepackage[mathlines]{lineno}          % 连续行号
\modulolinenumbers[5]                  % 每 5 行显示一次行号
\usepackage{fancyhdr}                   % 页眉与页脚
\usepackage[hidelinks]{hyperref}        % 可点击的引用和链接
\graphicspath{{figures/}}               % 图片放在 figures 文件夹

% 页面尺寸与章节编号，一般不需要修改。
\setlength{\textwidth}{16.5cm}
\setlength{\textheight}{22.6cm}
\setlength{\oddsidemargin}{-1mm}
\setlength{\evensidemargin}{-2mm}
\setlength{\topmargin}{-1cm}
\setlength{\headheight}{15pt}
\setlength{\headwidth}{\textwidth}
\setlength{\parindent}{0pt}             % 段落首行不缩进
\setlength{\parskip}{1.3ex}
\renewcommand{\thesection}{\Roman{section}.}
\renewcommand{\thesubsection}{\thesection\Alph{subsection}.}
\renewcommand{\thesubsubsection}{\thesubsection\arabic{subsubsection}.}
\makeatletter
\renewcommand\@biblabel[1]{$^{#1}$}
\makeatother
\setlength{\bibhang}{0em}


\setlength{\emergencystretch}{3em}

% 图表、表注和后续算法统一使用 8 pt、单倍行距。
\newcommand{\compactformat}{\setstretch{1}\fontsize{8}{9.6}\selectfont}
\makeatletter
\long\def\@makecaption#1#2{%
  \vskip\abovecaptionskip
  \begingroup
    \compactformat
    \sbox\@tempboxa{#1: #2}%
    \ifdim\wd\@tempboxa>\hsize
      #1: #2\par
    \else
      \global\@minipagefalse
      \hb@xt@\hsize{\hfil\box\@tempboxa\hfil}%
    \fi
  \endgroup
  \vskip\belowcaptionskip}
\makeatother

% 优先在段落边界分页，避免孤行及从断词处跨页；长段落仍允许分页。
\clubpenalty=10000
\widowpenalty=10000
\displaywidowpenalty=10000
\brokenpenalty=10000
\interlinepenalty=150

% 所有页面右上角显示页码，页眉为细横线，页脚留空且无线条。
\pagestyle{fancy}
\fancyhf{}
\fancyhead[R]{\thepage}
\renewcommand{\headrulewidth}{0.4pt}
\renewcommand{\footrulewidth}{0pt}
\fancypagestyle{plain}{%
  \fancyhf{}
  \fancyhead[R]{\thepage}
  \renewcommand{\headrulewidth}{0.4pt}
  \renewcommand{\footrulewidth}{0pt}
}


% 全文单栏；图表自动适配页面宽度。

% Allow nearby text and floats to share a page; avoid sparse float-only pages.
\renewcommand{\topfraction}{0.95}
\renewcommand{\bottomfraction}{0.90}
\renewcommand{\textfraction}{0.05}
\renewcommand{\floatpagefraction}{0.98}
\setcounter{topnumber}{3}
\setcounter{bottomnumber}{3}
\setcounter{totalnumber}{6}
\setlength{\textfloatsep}{12pt plus 2pt minus 2pt}
\setlength{\intextsep}{12pt plus 2pt minus 2pt}


% Algorithm title, pseudocode and note form one indivisible 8 pt block.
\newcounter{algorithm}
\makeatletter
\newcommand{\fnum@algorithm}{Algorithm \thealgorithm}
\newenvironment{algorithmblock}{%
  \begin{table}[H]\begin{minipage}{\linewidth}%
  \compactformat\let\footnotesize\compactformat
  \def\@captype{algorithm}%
}{\end{minipage}\end{table}}
\makeatother

\makeatletter\providecommand{\ext@algorithm}{loa}\makeatother

\setlength{\floatsep}{8pt plus 1pt minus 1pt}

\begin{document}

\raggedbottom

\begin{center}

{\Large\bfseries Integrated Optimization of Catheter Number and Spatial Configuration for Individualized Template-Guided High-Dose-Rate Interstitial Brachytherapy\par}

\vspace{6mm}

Qinglong~Yao, B.S.\textsuperscript{1}, Bo~Liu, Ph.D.\textsuperscript{1, 2,} *, Zhuo~Xiao, M.S.\textsuperscript{1}, Haitao~Sun, M.S.\textsuperscript{3}, Fugen~Zhou, Ph.D.\textsuperscript{1}, Junjie~Wang, M.D., Ph.D.\textsuperscript{3}, Qiuwen~Wu, Ph.D.\textsuperscript{4}, Ping~Jiang, M.D.\textsuperscript{3,} *\par

\vspace{4mm}

{\small

\emph{\textsuperscript{1} School of Astronautics, Beihang University, Beijing 100191, People's Republic of China}\par

\emph{\textsuperscript{2} State Key Laboratory of High-Efficiency Reusable Aerospace Transportation Technology, Beijing 102206, PR China}\par

\emph{\textsuperscript{3} Department of Radiation Oncology, Peking University Third Hospital, Beijing 100191, People's Republic of China}\par

\emph{\textsuperscript{4} Department of Radiation Oncology, Duke University Medical Center, Durham, North Carolina 27708, USA}\par

}

\end{center}

\textbf{Corresponding Authors}\par

\begingroup\small\setlength{\parskip}{0.4ex}

Bo Liu, Ph.D.\par

Image Processing Center\par

Beihang University\par

Beijing, 100191\par

P.R. China\par

Email: \href{mailto:bo.liu@buaa.edu.cn}{bo.liu@buaa.edu.cn}\par

\smallskip

Ping Jiang, M.D.\par

Department of Radiation Oncology\par

Peking University Third Hospital\par

Beijing, 100191\par

P.R. China\par

Email: \href{mailto:Jiangping@bjmu.edu.cn}{Jiangping@bjmu.edu.cn}\par

\smallskip

\endgroup

\clearpage

\begingroup\setlength{\parskip}{0.5ex}

\begin{abstract}

\noindent \textbf{Background}: Individualized template-guided high-dose-rate interstitial brachytherapy (HDR-ISBT) permits flexible catheter placement, but catheter number and configuration remain difficult to determine automatically.\par

\noindent \textbf{Purpose}: To develop and retrospectively evaluate an automated planning approach integrating catheter number determination with spatial configuration and dwell-time optimization for individualized template-guided HDR-ISBT.\par

\noindent \textbf{Methods}: We retrospectively replanned 94 fractions from 26 patients with locally recurrent cervical cancer. Feasible straight-catheter trajectories were generated, and an initial catheter pool was constructed through dose-guided selection and random supplementation. Group-sparsity optimization using the fast iterative shrinkage-thresholding algorithm (FISTA), followed by sequential pruning, determined catheter number without prespecifying the final count. Simulated annealing (SA)-based catheter exchange and local geometric fine-tuning subsequently optimized spatial configuration at that count, with dwell-time reoptimization. Automated plans were compared with clinical plans and across optimization stages. Ablation analyses assessed the azimuthal constraint, initial-pool construction, and sparsity weighting. Dosimetric metrics were evaluated after normalization to a high-risk clinical target volume (HR-CTV) D\textsubscript{90} of 600 cGy. Statistical inference used patient-level two-sided Wilcoxon signed-rank tests with a significance threshold of \emph{P} < 0.01.\par

\noindent \textbf{Results}: Compared with clinical plans, automated plans reduced the median catheter count from six to four (\emph{P} < 0.001), with fewer catheters in 78 of 94 plans. HR-CTV D\textsubscript{95} showed a median paired decrease of 5.7 cGy (\emph{P} = 0.002), whereas V\textsubscript{150} and V\textsubscript{200} did not differ significantly. Colon D\textsubscript{2cc} and urethral D\textsubscript{0.1cc} showed median paired decreases of 3.3 cGy (\emph{P} = 0.008) and 15.7 cGy (\emph{P} < 0.001), respectively; other organ-at-risk dose differences were not significant. Stage-wise analysis showed that Group-Sparsity optimization yielded a median of five catheters versus six clinically (\emph{P} = 0.016), with a small D\textsubscript{95} decrease and lower urethral D\textsubscript{0.1cc}. Pruning further reduced the median count to four but decreased D\textsubscript{95} and increased V\textsubscript{150}, V\textsubscript{200}, and organ-at-risk doses. Subsequent Catheter Exchange improved D\textsubscript{95} and reduced V\textsubscript{150}, V\textsubscript{200}, and several organ-at-risk doses at the retained count. Ablation analyses showed that differential weighting yielded fewer catheters than uniform weighting within the same dose-guided pool, accompanied by higher V\textsubscript{200} and colon D\textsubscript{2cc}. Under uniform weighting, dose-guided and random pools showed no significant differences in catheter count or dosimetry. Median total optimization time was 510.6 s.\par

\noindent \textbf{Conclusions}: The integrated automated approach reduced catheter use with a small HR-CTV D\textsubscript{95} trade-off, while subsequent Catheter Exchange improved dosimetric performance at the retained count. Prospective pre-implantation planning and physical validation are required before clinical implementation.\par

\end{abstract}

{\small\noindent\textbf{Keywords:} High-dose-rate interstitial brachytherapy; individualized template; automated treatment planning; catheter configuration optimization; group sparsity; simulated annealing\par}

\endgroup

\clearpage

\setstretch{1.5}

\linenumbers

\section{Introduction}\label{sec:1}

High-dose-rate interstitial brachytherapy (HDR-ISBT) is an important salvage treatment option for patients with locally recurrent cervical cancer after pelvic irradiation.\cite{ref1,ref2,ref3,ref4} In HDR-ISBT, achieving the intended dose distribution requires appropriate planning of the catheter configuration (including catheter number, positions, and orientations) and dwell times, together with accurate reproduction of the planned spatial configuration during implantation.\cite{ref5,ref6,ref7,ref8,ref9,ref10,ref11,ref12,ref13} In recent years, three-dimensional (3D) printed individualized templates (3D-PITs) have been developed to guide catheter insertion along predefined trajectories.\cite{ref14,ref15,ref16,ref17,ref18,ref19,ref20,ref21,ref22,ref23,ref24} However, clinical treatment planning systems offer limited dedicated support for designing template-guided catheter configurations before implantation. Consequently, determining the catheter configuration still relies heavily on manual trial-and-error and planner experience. In particular, the difficulty of systematically evaluating dosimetric trade-offs across different catheter counts may lead to conservative use of additional catheters.\cite{ref9,ref10,ref11,ref12,ref25,ref26} Furthermore, manually assessing three-dimensional collisions between nonparallel catheters is challenging, and planners often favor near-parallel configurations, potentially limiting the geometric flexibility afforded by individualized templates. Some of these configurations even include catheter trajectories outside the transvaginal template.\cite{ref25} Therefore, an integrated automated planning framework for 3D-PIT-guided HDR-ISBT is needed to determine catheter number and optimize catheter spatial configuration and dwell times under anatomical and geometric constraints, with the aim of reducing unnecessary catheter use while balancing target coverage and organ-at-risk sparing.

Previous studies have explored various aspects of automated HDR brachytherapy planning. Inverse planning models have been widely studied,\cite{ref27,ref28,ref29,ref30,ref31,ref32,ref33,ref34,ref35,ref36,ref37} and one common approach is to optimize dwell times with the catheter number, spatial configuration, and candidate dwell positions specified in advance.\cite{ref38,ref39,ref40,ref41,ref42,ref43,ref44,ref45,ref46} Other studies have extended optimization to catheter spatial configuration,\cite{ref26,ref47,ref48,ref49,ref50,ref51,ref52,ref53,ref54} including joint optimization of catheter positions and dwell times.\cite{ref55,ref56} Template-specific frameworks have also integrated catheter design with dose optimization.\cite{ref25,ref57} Deep learning-based multi-atlas registration has been used to predict catheter configurations before dwell-time optimization.\cite{ref58} Research on catheter spatial configuration has also extended to curved trajectories,\cite{ref59,ref60,ref61,ref62,ref63,ref64} with related work on instrument mechanics in curved channels\cite{ref65} and automated channel design for patient-specific applicators.\cite{ref66} Determining catheter number is another key planning decision.\cite{ref67} One approach uses geometric-coverage-based greedy selection without a prespecified final count, followed by dwell-time optimization and plan selection based on target coverage and catheter count.\cite{ref26} Other approaches optimize plans at several prespecified counts and compare their dosimetric performance to assess the feasibility and trade-offs of catheter reduction.\cite{ref25,ref48,ref55,ref68} Catheter selection has also been incorporated into inverse planning through compressed sensing-inspired heuristics\cite{ref69} and group-sparsity-based optimization.\cite{ref70,ref71} By penalizing the dwell-time vector of each catheter as a group, group-sparsity optimization can drive all dwell times on a catheter to zero, enabling catheter selection and dwell-time optimization within the same problem. The fast iterative shrinkage-thresholding algorithm (FISTA)\cite{ref72} has been used to solve such group-sparsity problems.\cite{ref70,ref71} Sparse catheter selection has also been combined with multicriteria dose optimization to generate candidate plans with different catheter counts, spatial configurations, and dosimetric trade-offs.\cite{ref73}

Together, these studies provide a solid methodological foundation for automated catheter planning. However, the scope of optimization depends on which planning variables are fixed in advance. When catheter geometry and candidate dwell positions are fixed, dose optimization is limited to adjusting dwell times.\cite{ref27,ref32} Spatial optimization at a prespecified catheter count allows catheter geometry to vary, but the count remains an input to each optimization run.\cite{ref25,ref55} Comparing plans across several counts can guide the choice of a final count, which may still depend on clinical judgment.\cite{ref74} Group-sparsity optimization can determine the catheter count for a given plan,\cite{ref70,ref73} but selecting from predefined candidate trajectories leaves the geometry of each retained trajectory unchanged.\cite{ref70,ref71} Optimizing catheter geometry and reoptimizing dwell times at the selected count can therefore provide additional flexibility to manage trade-offs among target coverage, high-dose volumes, and organ-at-risk doses.\cite{ref50,ref55,ref68} These considerations motivate an integrated workflow for 3D-PIT-guided HDR-ISBT that links template-constrained candidate generation with adaptive catheter count determination and subsequent optimization of catheter geometry and dwell times.

We developed and retrospectively evaluated an integrated automated planning framework for straight-catheter HDR-ISBT in patients with locally recurrent cervical cancer, based on a protocol using a combined transvaginal and transperineal 3D-PIT.\cite{ref20,ref21} The framework generates feasible candidate catheter trajectories, constructs an initial pool, and uses group-sparsity optimization and sequential pruning to select catheters without prespecifying the final count. At the resulting catheter count, the framework refines the spatial configuration through catheter exchange and local geometric fine-tuning, with dwell times reoptimized during configuration updates. We compared automated plans with clinical plans for 94 fractions from 26 patients to assess catheter use, target coverage, and organ-at-risk sparing. We also measured computation time, compared catheter use and dosimetry across stages, and used ablation analyses to assess the effects of selected design choices.

\section{Materials and Methods}\label{sec:2}

\subsection{Patient Cohort}\label{sec:2-1}

This retrospective study included 94 fraction plans from 26 patients with locally recurrent cervical cancer treated at Peking University Third Hospital between November 2017 and April 2023, all of whom had received prior pelvic external beam radiotherapy. Patients underwent 3D-PIT-guided HDR-ISBT at 6 Gy per fraction. Eligibility required that, after normalization of each clinical plan to a D\textsubscript{90} of 600 cGy for the high-risk clinical target volume (HR-CTV), all applicable dose limits for organs at risk (OARs) were satisfied: bladder D\textsubscript{2cc} \ensuremath{\leq} 550 cGy; rectum, colon, and small-intestine D\textsubscript{2cc} \ensuremath{\leq} 430 cGy; and urethral D\textsubscript{0.1cc} \ensuremath{\leq} 550 cGy. These institutional dose limits were informed by EMBRACE II.\cite{ref75} OARs without available contours were excluded from both eligibility assessment and automated optimization. The median number of included fractions per patient was 3 (interquartile range [IQR], 3--5; range, 1--7). The median HR-CTV was 53 cm\textsuperscript{3} (IQR, 33--73; range, 14--247). The median number of clinical catheters per fraction was 6 (IQR, 5--7; range, 3--11). The study was approved by the Ethics Committee of Peking University Third Hospital (approval No. IRB00006761-M2019210).

Computed tomography (CT) simulation and structure delineation were performed for each fraction. Post-insertion CT scans were acquired with the template and catheters in place (Brilliance Big Bore, Philips Healthcare, Cleveland, OH, USA; 120 kV, 150 mA, 3-mm slices). The gross tumor volume (GTV) was delineated based on CT-visible disease, with reference to diagnostic imaging and clinical examination. The HR-CTV included the GTV plus a margin of 0.3--0.5 cm. In cases with uterine corpus or vaginal involvement, the HR-CTV also included 3/4 of the uterine corpus or the entire vagina, thereby expanding the target volume to encompass surrounding tissues. Due to anatomical variability, not all OARs were considered clinically relevant for dosimetric evaluation and contoured in every fraction.

Clinical plans were generated using Eclipse 15.5 (Varian Medical Systems, USA) with volume-based optimization and dose shaper tools. All plans were prepared by a medical physicist with more than 10 years of experience following institutional protocols based on EMBRACE II\cite{ref75} and reviewed and approved by a senior radiation oncologist with more than 15 years of experience before treatment.

\subsection{Automated Catheter Planning Framework}\label{sec:2-2}

The proposed framework integrates catheter generation, catheter number determination, and spatial configuration and dwell-time optimization for 3D-PIT-guided HDR-ISBT using straight catheters. Given patient anatomical contours, template geometry, and dosimetric objectives, planning proceeds through three sequential stages, as illustrated in Figure~\ref{fig:1}.

\FloatBarrier
\begin{figure}[H]
\centering
\includegraphics[width=0.92\linewidth,height=0.78\textheight,keepaspectratio]{figures20261114/figure1.jpg}
\caption{Overview of the integrated automated planning workflow for 3D-PIT-guided HDR-ISBT. The input insets show two views of the patient's anatomical contours and the transvaginal template, with the HR-CTV shown in green. The output insets show an example catheter configuration with markers indicating source dwell positions and a three-dimensional view of the anatomical contours, catheter trajectories, and combined template with a patient-specific transperineal plate and guide channels. HR-CTV, high-risk clinical target volume; FISTA, fast iterative shrinkage-thresholding algorithm; SA, simulated annealing.}\label{fig:1}
\end{figure}

\begin{enumerate}
\item Candidate trajectory generation and initial-pool construction. Candidate trajectories satisfying anatomical and geometric constraints are generated. Dose-guided selection and random supplementation then construct a geometrically compatible initial catheter pool for subsequent group-sparsity-based catheter selection and dwell-time optimization.
\item Adaptive catheter number determination. Group-sparsity optimization and sequential pruning select the retained catheters and determine their number without prespecifying the final catheter count. Dwell-time reoptimization is used to assess target-coverage feasibility during catheter reduction.
\item Catheter spatial configuration optimization. At the determined catheter count, catheter exchange and local geometric fine-tuning adjust the spatial distribution of available dwell positions. Dwell times are reoptimized to adapt the dose distribution to each updated configuration.
\end{enumerate}

Dwell-time optimization during catheter number determination and spatial configuration optimization uses a common dose term, defined as follows:

\begin{equation}\label{eq:1}
\begin{aligned}F_{\text{dose}}\left( t \right) &= \frac{\lambda_{T}}{V_{T}}\sum_{i \in \mathcal{A}_{T}\left( t \right)}^{}{v_{i}^{T}\left\lbrack D_{T}^{\text{LB}} - d_{T,i}(t) \right\rbrack^{2}}\\ &{}+\sum_{r \in \mathcal{O}}^{}{\frac{\lambda_{r}}{V_{r}}\sum_{i \in \mathcal{A}_{r}\left( t \right)}^{}{v_{\text{ri}}^{O}\left\lbrack d_{r,i}\left( t \right) - D_{r}^{\text{UB}} \right\rbrack^{2}}}.\end{aligned}
\end{equation}

Here, $t$ is the dwell-time vector, $T$ the HR-CTV, $\mathcal{O}$ the set of evaluated OARs, and $d$ the doses at sampling points in voxels overlapping the contoured structures. At sampling point $i$, $v_{i}^{T}$ and $v_{\text{ri}}^{O}$ are the HR-CTV and OAR $r$ volume fractions, respectively; $V_{T}$ and $V_{r}$ are their sums over all sampling points in the respective structures. $D_{T}^{\text{LB}}$ denotes the lower bound for HR-CTV D\textsubscript{90}, whereas $D_{r}^{\text{UB}}$ denotes the upper bound for OAR $r$: D\textsubscript{2cc} for the bladder, rectum, colon, and small intestine, and D\textsubscript{0.1cc} for the urethra. The set $\mathcal{A}_{T}\left( t \right)$ includes target points at or above the current D\textsubscript{90} but below $D_{T}^{\text{LB}}$; $\mathcal{A}_{r}\left( t \right)$ includes OAR points above $D_{r}^{\text{UB}}$ but at or below the current corresponding dose-volume metric. The dose weights were $\lambda_{T}$ = $\lambda_{r}$ = 100.

For each fraction, $D_{T}^{\text{LB}}$ and $D_{r}^{\text{UB}}$ were set to the corresponding dose-volume metrics of the original, unrescaled clinical plan. These bounds remained fixed throughout optimization and feasibility assessment. Normalization to HR-CTV D\textsubscript{90} = 600 cGy was used only for eligibility assessment and plan comparisons.

The common dose term $F_{\text{dose}}\left( t \right)$ is augmented by the group-sparsity penalty $F_{\text{sparse}}\left( t \right)$ during Group-Sparsity and by the dwell-time smoothness penalty $F_{\text{smooth}}\left( t \right)$ during spatial configuration optimization at a fixed catheter count. Dwell-time reoptimization during coverage checks and Pruning uses $F_{\text{dose}}\left( t \right)$ alone.

\subsubsection{Candidate Trajectory Generation and Initial Pool Construction}\label{sec:2-2-1}

As illustrated in Figure~\ref{fig:2}(A), the catheter-entry region comprised the outer lateral surface of the proximal 15-mm segment of the transvaginal template and the template's proximal end face. Following the sampling approach of Yao et al.,\cite{ref25} the entry and target sampling planes were perpendicular to the patient's superior--inferior axis and were each offset outward by 5 mm from the inferior boundary of the transvaginal template and the superior boundary of the HR-CTV, respectively. On each plane, the sampling region was defined as the convex hull of the intersections of that plane with lines connecting points in the catheter-entry region to HR-CTV contour points. Each straight candidate catheter trajectory was formed by connecting a pair of points, one sampled uniformly at random from each convex hull.

\FloatBarrier
\begin{figure}[H]
\centering
\includegraphics[width=0.92\linewidth,height=0.78\textheight,keepaspectratio]{figures20261114/figure2.jpg}
\caption{Candidate trajectory generation and initial-pool construction. (A) Generation and feasibility screening of candidate trajectories. $C_{e}$ and $C_{t}$ denote the area centroids of the convex hulls defining the entry and target sampling regions, and $P_{e}$ and $P_{t}$ denote points sampled on the entry and target sampling planes, respectively. The blue line connecting $P_{e}$ and $P_{t}$ illustrates a candidate trajectory. The azimuths $\theta_{e}$ and $\theta_{t}$ are measured from the positive \emph{x}-axis to the vector from the corresponding centroid to the sampled point in each plane. (B) Initial-pool construction using dose-guided core selection and random supplementation. Green and black trajectories in (B) represent core and randomly supplemented catheters, respectively.}\label{fig:2}
\end{figure}

Each trajectory was truncated at its distal exit point from the HR-CTV. The available space for candidate trajectory generation within the cylindrical template was defined as an interior region obtained by offsetting its lateral surface inward by 2 mm. This offset accounted for the radius of the 2-mm guide channel used for interstitial catheter insertion and a 1-mm outer-wall allowance adopted in the template design. A trajectory was retained if it intersected the boundaries of both the HR-CTV and individualized template at least twice, traversed at least 10 mm within the HR-CTV and 50 mm within this interior region, and did not intersect any contoured OAR along the retained trajectory.

To discourage pronounced catheter crossings, we designed an entry--target azimuthal consistency constraint for candidate trajectory generation. This constraint limits the difference between the azimuthal angles of the vectors from the respective convex-hull centroids to the sampled entry and target points, as illustrated in Figure~\ref{fig:2}(A). The entry and target azimuthal angles were calculated as $\theta_{e} = \operatorname{atan2}\left( y_{e} - {\overline{y}}_{e},\ x_{e} - {\overline{x}}_{e} \right)$ and $\theta_{t} = \operatorname{atan2}\left( y_{t} - {\overline{y}}_{t},\ x_{t} - {\overline{x}}_{t} \right)$, respectively. Here, $P_{e} = \left( x_{e},y_{e}\  \right)$ and $P_{t} = \left( x_{t},y_{t}\  \right)$ are the sampled entry point and target point, and $C_{e} = \left( {\overline{x}}_{e},{\overline{y}}_{e}\  \right)$ and $C_{t} = \left( {\overline{x}}_{t},{\overline{y}}_{t}\  \right)$ are the area centroids of the convex hulls defining the entry and target sampling regions. The minimum angular difference was constrained by Equation~\eqref{eq:2}:

\begin{equation}\label{eq:2}
\mathrm{\Delta}\theta = \min\left( \left| \theta_{e} - \theta_{t} \right|,\ 2\pi - \left| \theta_{e} - \theta_{t} \right| \right) \leq \frac{\pi}{6}.
\end{equation}

For each fraction, candidate generation terminated when either 3,000 feasible trajectories had been obtained or 300,000 sampling attempts had been reached.

To provide a dose-informed initial pool while retaining alternative trajectories for subsequent optimization, dose-guided core-catheter selection was combined with random supplementation, as shown in Figure~\ref{fig:2}(B). Random supplementation added geometrically compatible candidates not selected by the dose-guided procedure, reducing reliance on the proxy ranking during subsequent catheter number optimization. For dose-guided core selection, each candidate catheter $c$ was assigned a proxy-dose-based score reflecting target coverage and potential OAR irradiation, as defined in Equation~\eqref{eq:3}:

\begin{equation}\label{eq:3}
S_{c} = \sum_{i}^{}(w_{i}^{\text{cold}} - w_{i}^{\text{risk}})q_{\text{ci}},
\end{equation}

where $w_{i}^{\text{cold}}$ and $w_{i}^{\text{risk}}$ denote the target-coverage priority and OAR-risk weights at sampling point $i$, respectively, and $q_{\text{ci}} \in \left\lbrack 0,1 \right\rbrack$ is the normalized proxy dose contribution from candidate catheter $c$ at that point. The proxy dose was calculated using dwell positions spaced 5 mm apart along the candidate catheter between its HR-CTV entry and exit points. The same 5-mm dwell-position spacing was used in all subsequent optimization steps.

At each iteration, the highest-scoring candidate with $S_{c} > 0$ was selected among those satisfying a minimum 2-mm separation from all previously selected catheters. The spacing constraint was imposed to prevent overlap between guide channels, with separation measured as the minimum three-dimensional Euclidean distance between the finite catheter centerline segments. Following each selection, the target-coverage priority weights were reduced according to the normalized proxy dose contribution from the newly selected catheter, and the OAR-risk weights were updated based on the cumulative proxy dose from the selected core catheters. Candidate scores were then recalculated to account for the updated target-coverage priorities and OAR-risk penalties. Core selection terminated when the sum of the target-coverage priority weights fell below 10\textsuperscript{-4} of its initial value, the 100-catheter pool limit was reached, or no geometrically eligible candidate with a positive score remained. The proxy dose calculation, weight initialization, and weight updates are detailed in Appendix A.1.

Following core selection, additional candidate catheters with valid dwell positions were drawn in random order from the remaining feasible trajectories, subject to the same minimum 2-mm three-dimensional inter-catheter separation constraint. Random supplementation continued until the initial catheter pool contained 100 catheters or no eligible candidates remained.

\subsubsection{Adaptive Catheter Number Determination}\label{sec:2-2-2}

Group-sparsity optimization was performed using FISTA\cite{ref70,ref71} to jointly select catheters and optimize dwell times without prespecifying the final catheter count. The optimization problem was formulated as

\begin{equation}\label{eq:4}
\min_{t \geq 0}\left\lbrack F_{\text{dose}}\left( t \right) + F_{\text{sparse}}\left( t \right) \right\rbrack,\ with\ F_{\text{sparse}}\left( t \right) = \sum_{c = 1}^{C}{\omega_{c}(K)\left\| t_{c} \right\|_{2}^{\frac{1}{2}}}.
\end{equation}

Here, $C$ is the number of candidate catheters in the initial catheter pool, and $F_{\text{dose}}\left( t \right)$ is defined in Equation~\eqref{eq:1}. The term $F_{\text{sparse}}\left( t \right)$ is a group-half penalty on the dwell-time vector $t_{c}$ of each catheter, encouraging all dwell times of individual catheters to become zero and thereby excluding those catheters from the plan, as illustrated in Figure~\ref{fig:3}(A).\cite{ref71} Building on this formulation, we introduced a catheter-category multiplier $m_{c}$ and redesigned the update mechanism for the global sparsity strength $K$ to make it dose dependent. The catheter-specific sparsity weight was defined as

\FloatBarrier
\begin{figure}[H]
\centering
\includegraphics[width=0.92\linewidth,height=0.78\textheight,keepaspectratio]{figures20261114/figure3.jpg}
\caption{Adaptive catheter number determination using group-sparsity optimization and sequential pruning. In (B), each tentative removal was subject to the 5\% cumulative dwell-time limit and was accepted only if, after dwell-time reoptimization, the dose objective was finite and HR-CTV D\textsubscript{90} was at least 95\% of its value in the original, unrescaled clinical plan.}\label{fig:3}
\end{figure}

\begin{equation}\label{eq:5}
\omega_{c}\left( K \right) = Km_{c}\rho_{c},
\end{equation}

where $\rho_{c}$ is the normalization factor adopted from the group-sparsity formulation to account for differences in catheter dose loading and dwell-position count.\cite{ref71} The newly introduced multiplier $m_{c}$ assigns different relative sparsity penalties to the catheter categories established during initial-pool construction. We set $m_{c}$ = 0.7 for dose-guided core catheters and $m_{c}$ = 1.5 for randomly supplemented catheters. The lower multiplier for core catheters was intended to favor their retention while keeping all catheters eligible for removal.

To adapt catheter-removal pressure to dosimetric performance, the global sparsity strength was initialized at $K_{0}$ = 20 and increased between blocks typically comprising 500 iterations:

\begin{equation}\label{eq:6}
K \leftarrow K\left( 1 + \frac{0.2}{1 + 5\varepsilon_{\max}} \right).
\end{equation}

Here, $\varepsilon_{\max}$ is the maximum relative violation among active HR-CTV and OAR dose objectives with positive weights and valid dose-volume histogram (DVH) metrics. For each objective, the relative violation was calculated by dividing the amount of the violation by the corresponding dose bound and was set to zero when the bound was satisfied. Larger violations slowed the increase in sparsity strength.

To safeguard HR-CTV coverage, we incorporated a feasibility check into the Group-Sparsity stage. At each check, dwell times were reoptimized by minimizing $F_{\text{dose}}\left( t \right)$ using the bound-constrained limited-memory Broyden--Fletcher--Goldfarb--Shanno (L-BFGS-B) algorithm,\cite{ref76,ref77} with nonnegative bounds and a limit of 50 iterations. The first coverage check was performed at iteration 551. Thereafter, checks were triggered whenever the active catheter count fell below that recorded at the latest successful check, without waiting for a block boundary. A configuration passed if the reoptimized HR-CTV D\textsubscript{90} was at least 95\% of the corresponding original, unrescaled clinical-plan D\textsubscript{90} ($D_{T}^{\text{LB}}$). The 95\% threshold was selected empirically to limit the allowable decrease in HR-CTV D\textsubscript{90} relative to the clinical reference at each feasibility check.

Successful checks saved the active set and FISTA dwell times as the latest feasible anchor; L-BFGS-B dwell times were used only for assessment. Failure of the first coverage check terminated the workflow without proceeding to Pruning or spatial configuration optimization. A subsequent failure restored the latest feasible anchor and ended FISTA. FISTA also stopped when three consecutive eligible block-end comparisons showed no change in the active catheter set (the active-set stability criterion) or when the iteration count reached 20,000. The eligibility and counting rules for these comparisons are detailed in Appendix A.2.

Following the Group-Sparsity stage, sequential pruning was performed to explore further catheter reduction, as shown in Figure~\ref{fig:3}(B). Catheters were ranked once in ascending order of total dwell time in the final FISTA solution. The cumulative dwell time of removed catheters was capped at 5\% of that solution's total dwell time, with the ranking and cap fixed throughout Pruning. The same L-BFGS-B procedure was used to minimize $F_{\text{dose}}\left( t \right)$ for the complete active set before Pruning and after each tentative removal. Removal was accepted only if the reoptimized value of $F_{\text{dose}}\left( t \right)$ was finite and the reoptimized HR-CTV D\textsubscript{90} was at least 95\% of $D_{T}^{\text{LB}}$; otherwise, the last accepted solution was restored and Pruning stopped. This empirically selected D\textsubscript{90} tolerance allowed further catheter reduction while limiting the allowable decrease in reoptimized HR-CTV D\textsubscript{90} relative to the original, unrescaled clinical plan. Pruning also stopped if removing the next catheter would exceed the 5\% cumulative dwell-time limit or if only one catheter remained. The control flow for Group-Sparsity and Pruning is summarized in Algorithm A1 in Appendix A.2.

\subsubsection{Catheter Spatial Configuration Optimization}\label{sec:2-2-3}

Starting from the configuration retained after Pruning, spatial configuration optimization sought to improve the dose distribution while maintaining the determined catheter count. It comprised simulated annealing (SA)-based catheter exchange followed by local geometric fine-tuning. Once the catheter count was fixed, a dwell-time smoothness penalty was introduced to discourage abrupt changes between adjacent dwell times and promote more regular dwell-time distributions. Both steps minimized $F_{\text{dose}}\left( t \right) + F_{\text{smooth}}\left( t \right)$, with the smoothness penalty defined as:

\begin{equation}\label{eq:7}
F_{\text{smooth}}\left( t \right) = \frac{\lambda_{\text{smooth}}}{2}\sum_{c = 1}^{C}{\sum_{j = 1}^{n_{c} - 1}\left( t_{c,j + 1} - t_{c,j} \right)^{2}}.
\end{equation}

Here, $C$ is the retained catheter count, $t_{c,j}$ denotes the dwell time at position $j$ of catheter $c$, $n_{c}$ is its number of dwell positions, and $\lambda_{\text{smooth}}$ = 1.

Both steps used SA to search over catheter configurations and L-BFGS-B to minimize $F_{\text{dose}}\left( t \right) + F_{\text{smooth}}\left( t \right)$ over dwell times for each fixed configuration. Before each step, L-BFGS-B reoptimized dwell times to establish a baseline. Each proposed configuration was first checked for geometric feasibility. For feasible proposals, L-BFGS-B reoptimized dwell times for all catheters before objective evaluation, allowing each configuration to be assessed with an adapted dwell-time distribution. Both steps used the Metropolis acceptance criterion,\cite{ref78} allowing probabilistic acceptance of objective increases, and retained their respective lowest-objective solutions.

To explore alternative catheter combinations across the full feasible candidate set, each iteration of Catheter Exchange proposed replacing one randomly selected catheter with a candidate not already present in the current configuration, as illustrated in Figure~\ref{fig:4}(A). A proposal required valid dwell positions within the HR-CTV and at least 2-mm separation from the remaining catheters. The initial temperature was $T_{0} = max(10^{- 6},0.3F_{0})$, where $F_{0}$ was the baseline value of $F_{\text{dose}}\left( t \right) + F_{\text{smooth}}\left( t \right)$ after dwell-time reoptimization, and the temperature was reduced by a factor of 0.95 after each replacement evaluation. Catheter Exchange stopped after at most 1,000 outer iterations or more than 50 consecutive Metropolis rejections following dwell-time reoptimization.

\FloatBarrier
\begin{figure}[H]
\centering
\includegraphics[width=0.92\linewidth,height=0.78\textheight,keepaspectratio]{figures20261114/figure4.jpg}
\caption{Catheter spatial configuration optimization using SA-based catheter exchange and local geometric fine-tuning. Gray and orange trajectories indicate the original and proposed paths, respectively. The shared optimization loop in (C) is applied separately in both stages, with the catheter count held constant.}\label{fig:4}
\end{figure}

Starting from the configuration retained after Catheter Exchange, Fine-Tuning explored neighboring configurations without restricting trajectories to the sampled candidate set, as shown in Figure~\ref{fig:4}(B). At each iteration, one catheter was randomly selected, and the \emph{x}- or \emph{y}-coordinate of either its entry or target endpoint was perturbed uniformly within \ensuremath{\pm}1/50 of the corresponding sampling-plane extent along that axis. After perturbation, the target-side endpoint was reset to the HR-CTV surface by truncating or extending the updated trajectory. Proposals violating the candidate-generation geometric constraints or the 2-mm inter-catheter separation were rejected. Dwell positions were regenerated for geometrically feasible trajectories, and proposals without valid HR-CTV dwell positions were also rejected. Fine-Tuning used an empirically selected initial temperature of 100 and the same cooling schedule as Catheter Exchange. Fine-Tuning stopped after at most 1,000 perturbation attempts or more than 50 consecutive Metropolis rejections following dwell-time reoptimization.

\subsection{Ablation Study Design}\label{sec:2-3}

Four configurations were evaluated to assess the effects of the azimuthal constraint, initial-pool construction, and sparsity weighting:

\begin{enumerate}
\item Reference. The Reference configuration used the 30\textdegree{} azimuthal constraint, dose-guided core selection with random supplementation, and sparsity-weight multipliers of 0.7 and 1.5 for core and supplementary catheters, respectively. The azimuthal threshold was specified based on clinical experience, whereas the sparsity-weight multipliers were selected based on preliminary testing.
\item No azimuthal constraint. The azimuthal constraint was omitted during candidate generation and Fine-Tuning checks. All other Reference settings were retained.
\item Uniform-weight. The Reference initial catheter pool was used with a sparsity-weight multiplier of 1.0 for all catheters.
\item Random-pool. The initial catheter pool was constructed from the same eligible candidate set in random order under the same geometric criteria and pool limit of 100 catheters, also using a sparsity-weight multiplier of 1.0 for all catheters.
\end{enumerate}

Ablation comparisons assessed the azimuthal constraint (Reference versus No azimuthal constraint), differential sparsity weighting (Reference versus Uniform-weight), initial-pool construction (Uniform-weight versus Random-pool), and the combined effects of pool construction and weighting (Reference versus Random-pool). All other data, objectives, and settings were held constant. Procedure-specific random seeds were fixed across experimental configurations.

\subsection{Plan Evaluation and Statistical Analysis}\label{sec:2-4}

The final automated plans were compared with the clinical plans to evaluate catheter reduction and dosimetric performance. Stage-wise comparisons assessed clinical versus Group-Sparsity, Group-Sparsity versus Pruning, Pruning versus Catheter Exchange, and Catheter Exchange versus Fine-Tuning. Dosimetric comparisons across the four automated optimization stages used dwell times optimized with the same objective, $F_{\text{dose}}\left( t \right) + F_{\text{smooth}}\left( t \right)$. For Group-Sparsity and Pruning, independent dwell-time reoptimizations were performed with the retained catheter configurations fixed; these evaluations did not affect subsequent optimization. Catheter Exchange and Fine-Tuning metrics were obtained from their respective retained best solutions.

Dosimetric metrics included HR-CTV D\textsubscript{95}, V\textsubscript{150}, and V\textsubscript{200}; bladder, rectum, colon, and small-intestine D\textsubscript{2cc}; and urethral D\textsubscript{0.1cc}. For comparison, all clinical plans and automated plans from the evaluated stages and configurations were individually normalized to HR-CTV D\textsubscript{90} = 600 cGy, with automated plans normalized after the relevant optimization or independent dwell-time reoptimization. Dosimetric metrics, total dwell time, and the dwell-time smoothness penalty were then calculated from the normalized plans. Additional metrics included candidate and initial-pool counts, catheter counts after Group-Sparsity and Pruning, and active dwell-position counts for Catheter Exchange and Fine-Tuning. Computational metrics included iteration counts and termination behavior for FISTA, Catheter Exchange, and Fine-Tuning, as well as candidate-generation, stage-specific, and total computation times. Total optimization time spanned candidate generation through Fine-Tuning. AAPM TG-43 unit-time dose rates\cite{ref79} were calculated dynamically for each configuration using an NVIDIA GeForce RTX 2080 Ti GPU, while optimization ran on an Intel Xeon E5-2620 v4 CPU.

Plan metrics were summarized at the fraction-plan level as medians with the first and third quartiles (Q1 and Q3). Stage-wise paired changes were calculated as the later stage minus the preceding stage. To account for multiple fractions per patient, valid fraction-level values were averaged within each patient and comparison condition, and the resulting patient-level means were compared using two-sided Wilcoxon signed-rank tests. Patients without valid paired data for a metric were excluded from that comparison. The nominal significance level was 0.05. Bonferroni correction across the five OAR metrics yielded a threshold of \emph{P} < 0.01, which was also applied to all other patient-level paired metrics.

\FloatBarrier
\section{Results}\label{sec:3}

\subsection{Clinical versus Automated Plans}\label{sec:3-1}

All 94 plans completed the automated workflow. Figure~\ref{fig:5} shows the paired catheter counts and their distributions for clinical and final automated plans. Of the 94 plans, 78 (83.0\%) used fewer catheters, 11 (11.7\%) were unchanged, and five (5.3\%) used more. Median catheter count decreased from six to four, with a median paired change of \ensuremath{-}2 (\emph{P} < 0.001).

\FloatBarrier
\begin{figure}[H]
\centering
\includegraphics[width=0.92\linewidth,height=0.78\textheight,keepaspectratio]{figures20261114/figure5.jpg}
\caption{Clinical versus final automated catheter counts. (A) Joint-frequency bubble plot. Circle area and the number inside indicate the frequency of each catheter-count combination; the dashed line indicates equal catheter counts. (B) Marginal-frequency distributions of clinical and final automated catheter counts. Bar heights and labels indicate the number of fraction plans at each catheter count.}\label{fig:5}
\end{figure}

After normalization to HR-CTV D\textsubscript{90} = 600 cGy, the final automated plans showed a median paired decrease of 5.7 cGy in HR-CTV D\textsubscript{95} (\emph{P} = 0.002), alongside lower colon D\textsubscript{2cc} and urethral D\textsubscript{0.1cc} (median paired changes, \ensuremath{-}3.3 and \ensuremath{-}15.7 cGy, respectively; both \emph{P} < 0.01). No significant differences were detected in HR-CTV V\textsubscript{150} or V\textsubscript{200}, the other OAR dose metrics, or total dwell time. Table~\ref{tab:1} summarizes the clinical and final automated plan metrics. Figure~\ref{fig:6} presents a representative fraction selected because its clinical and final automated catheter counts matched the respective cohort medians.

\FloatBarrier
\begin{figure}[H]
\centering
\includegraphics[width=0.92\linewidth,height=0.65\textheight,keepaspectratio]{figures20261114/figure6.jpg}
\caption{Comparison of clinical and final automated plans for a representative fraction. (A--C) Clinical plan and (D--F) final automated plan, showing axial dose distributions, sagittal dose distributions, and three-dimensional catheter configurations, respectively. (G) HR-CTV and (H) OAR cumulative dose--volume histograms comparing both plans. The HR-CTV volume was 56.44 cm\textsuperscript{3}, and catheter counts were six and four for the clinical and final automated plans, respectively. Both plans were individually normalized to an HR-CTV D\textsubscript{90} of 600 cGy.}\label{fig:6}
\end{figure}

\begin{table}[H]
\centering
\begin{minipage}{\linewidth}
\compactformat\let\footnotesize\compactformat\setlength{\parskip}{0pt}
\setlength{\tabcolsep}{3pt}
\renewcommand{\arraystretch}{1}
\caption{Clinical and final automated plan metrics.}\label{tab:1}
\begin{tabularx}{\linewidth}{@{}>{\raggedright\arraybackslash}p{0.18\linewidth}*{3}{>{\raggedright\arraybackslash}X}>{\raggedright\arraybackslash}p{0.07\linewidth}>{\raggedright\arraybackslash}p{0.10\linewidth}@{}}
\specialrule{0.5pt}{1.5pt}{1.5pt}
\textbf{\textbf{Metric}} & \textbf{\textbf{Clinical plans}} & \textbf{\textbf{Automated plans}} & \textbf{\textbf{Paired change}} & \textbf{\textbf{\emph{P} value}} & \textbf{\textbf{Patients/\newline fractions}} \\
\specialrule{1.5pt}{1.5pt}{1.5pt}
Number of catheters & 6 (5, 7) & 4 (3, 5) & \ensuremath{-}2 (\ensuremath{-}3, \ensuremath{-}1) & <0.001 & 26/94 \\
HR-CTV D\textsubscript{95}, cGy & 542.4 (533.1, 548.9) & 538.2 (527.9, 543.8) & \ensuremath{-}5.7 (\ensuremath{-}12.3, 2.2) & 0.002 & 26/94 \\
HR-CTV V\textsubscript{150}, \% & 55.0 (51.7, 57.2) & 55.9 (53.2, 58.0) & 1.3 (\ensuremath{-}1.8, 2.9) & 0.020 & 26/94 \\
HR-CTV V\textsubscript{200}, \% & 31.2 (26.5, 34.5) & 31.0 (25.2, 35.0) & 0.5 (\ensuremath{-}4.0, 4.1) & 0.568 & 26/94 \\
Bladder D\textsubscript{2cc}, cGy & 430.3 (351.3, 464.3) & 423.9 (325.4, 463.4) & \ensuremath{-}1.8 (\ensuremath{-}11.4, 2.8) & 0.096 & 26/94 \\
Rectum D\textsubscript{2cc}, cGy & 367.9 (329.2, 394.8) & 363.3 (325.6, 393.3) & \ensuremath{-}0.4 (\ensuremath{-}6.9, 4.2) & 0.058 & 26/94 \\
Colon D\textsubscript{2cc}, cGy & 185.3 (105.2, 323.8) & 183.4 (101.0, 314.9) & \ensuremath{-}3.3 (\ensuremath{-}9.0, 0.4) & 0.008 & 19/57 \\
Small intestine D\textsubscript{2cc}, cGy & 133.2 (103.5, 170.9) & 121.0 (102.3, 183.4) & \ensuremath{-}5.2 (\ensuremath{-}14.4, \ensuremath{-}0.7) & 0.248 & 11/19 \\
Urethra D\textsubscript{0.1cc}, cGy & 344.9 (194.7, 446.2) & 314.0 (186.3, 420.3) & \ensuremath{-}15.7 (\ensuremath{-}31.1, \ensuremath{-}5.8) & <0.001 & 25/90 \\
Total dwell time, s & 246.4 (184.3, 295.2) & 244.0 (177.2, 301.3) & 2.2 (\ensuremath{-}7.3, 14.1) & 0.107 & 26/94 \\
\specialrule{1.5pt}{1.5pt}{1.5pt}
\end{tabularx}
\par\vspace{3pt}
\noindent Note: Automated plans refer to the final plans retained after Fine-Tuning. Values are presented as medians (first quartile, third quartile), except for \emph{P} values and patient/fraction counts. Paired changes were calculated at the fraction-plan level as final automated minus clinical. \emph{P} values were obtained from two-sided Wilcoxon signed-rank tests on patient-level means, with significance assessed at \emph{P} < 0.01. HR-CTV, high-risk clinical target volume.\par
\end{minipage}
\end{table}

\subsection{Stage-Wise Results and Computation Time}\label{sec:3-2}

Candidate generation reached 3,000 feasible trajectories in all 94 plans. The median initial catheter pool contained 62 catheters (IQR, 56 to 68). Median counts were 2.5 core catheters (IQR, 1 to 3) and 60 supplementary catheters (IQR, 55 to 64).

The median catheter count after Group-Sparsity was five, compared with six in the clinical plans; the paired difference did not meet the prespecified significance threshold (\emph{P} = 0.016). Relative to the clinical plans, HR-CTV D\textsubscript{95} decreased by a median of 3.8 cGy (\emph{P} = 0.007), and urethral D\textsubscript{0.1cc} decreased by 10.3 cGy (\emph{P} < 0.001). Compared with Group-Sparsity, Pruning further reduced the median catheter count from five to four (\emph{P} < 0.001), accompanied by lower HR-CTV D\textsubscript{95} and higher V\textsubscript{150}, V\textsubscript{200}, all evaluated OAR dose metrics, and total dwell time (all \emph{P} < 0.01).

\begin{table}[H]
\centering
\begin{minipage}{\linewidth}
\compactformat\let\footnotesize\compactformat\setlength{\parskip}{0pt}
\setlength{\tabcolsep}{3pt}
\renewcommand{\arraystretch}{1}
\caption{Catheter counts, dosimetric indices, and dwell-time characteristics of clinical plans and automated plans after Group-Sparsity, Pruning, Catheter Exchange, and Fine-Tuning.}\label{tab:2}
\begin{tabularx}{\linewidth}{@{}>{\raggedright\arraybackslash}p{0.14\linewidth}@{\hspace{2pt}}*{5}{>{\raggedright\arraybackslash}X}@{}}
\specialrule{0.5pt}{1.5pt}{1.5pt}
\textbf{\textbf{Metric}} & \textbf{\textbf{Clinical}} & \textbf{\textbf{Group-Sparsity}} & \textbf{\textbf{Pruning}} & \textbf{\textbf{Catheter Exchange}} & \textbf{\textbf{Fine-Tuning}} \\
\specialrule{1.5pt}{1.5pt}{1.5pt}
Number of catheters & \mbox{6 (5, 7)} & \mbox{5 (4, 6)} & \mbox{4 (3, 5)} & \mbox{4 (3, 5)} & \mbox{4 (3, 5)} \\
HR-CTV D\textsubscript{95}, cGy & \mbox{542.4 (533.1, 548.9)} & \mbox{537.2 (528.9, 543.6)} & \mbox{533.9 (526.0, 540.5)} & \mbox{537.3 (529.9, 544.2)} & \mbox{538.2 (527.9, 543.8)} \\
HR-CTV V\textsubscript{150}, \% & \mbox{55.0 (51.7, 57.2)} & \mbox{55.1 (51.8, 57.3)} & \mbox{56.2 (53.9, 58.9)} & \mbox{55.5 (52.7, 57.8)} & \mbox{55.9 (53.2, 58.0)} \\
HR-CTV V\textsubscript{200}, \% & \mbox{31.2 (26.5, 34.5)} & \mbox{31.7 (25.2, 34.9)} & \mbox{33.0 (28.6, 36.9)} & \mbox{31.5 (25.0, 34.8)} & \mbox{31.0 (25.2, 35.0)} \\
Bladder D\textsubscript{2cc}, cGy & \mbox{430.3 (351.3, 464.3)} & \mbox{424.0 (324.7, 470.2)} & \mbox{429.5 (334.4, 472.9)} & \mbox{423.9 (330.8, 464.7)} & \mbox{423.9 (325.4, 463.4)} \\
Rectum D\textsubscript{2cc}, cGy & \mbox{367.9 (329.2, 394.8)} & \mbox{364.0 (324.1, 394.8)} & \mbox{373.0 (328.9, 405.3)} & \mbox{368.2 (325.6, 393.4)} & \mbox{363.3 (325.6, 393.3)} \\
Colon D\textsubscript{2cc}, cGy & \mbox{185.3 (105.2, 323.8)} & \mbox{186.4 (103.5, 313.6)} & \mbox{192.6 (104.9, 317.9)} & \mbox{183.5 (100.2, 316.0)} & \mbox{183.4 (101.0, 314.9)} \\
Small intestine D\textsubscript{2cc}, cGy & \mbox{133.2 (103.5, 170.9)} & \mbox{120.4 (101.4, 177.1)} & \mbox{124.7 (101.5, 186.2)} & \mbox{119.5 (102.8, 181.7)} & \mbox{121.0 (102.3, 183.4)} \\
Urethra D\textsubscript{0.1cc}, cGy & \mbox{344.9 (194.7, 446.2)} & \mbox{313.8 (183.8, 427.4)} & \mbox{319.0 (184.4, 429.6)} & \mbox{314.3 (187.2, 421.9)} & \mbox{314.0 (186.3, 420.3)} \\
Total dwell time, s & \mbox{246.4 (184.3, 295.2)} & \mbox{243.3 (178.7, 303.1)} & \mbox{247.3 (184.3, 310.8)} & \mbox{245.1 (179.9, 311.8)} & \mbox{244.0 (177.2, 301.3)} \\
Number of active dwell positions & \mbox{---} & \mbox{---} & \mbox{---} & \mbox{38.0 (26.3, 49.0)} & \mbox{37.0 (27.0, 49.8)} \\
Dwell-time smoothness penalty & \mbox{---} & \mbox{---} & \mbox{---} & \mbox{8.9 (3.0, 36.2)} & \mbox{5.8 (0.9, 14.4)} \\
\specialrule{1.5pt}{1.5pt}{1.5pt}
\end{tabularx}
\par\vspace{3pt}
\noindent Note: Values are presented as medians (first quartile, third quartile). HR-CTV, high-risk clinical target volume.\par
\end{minipage}
\end{table}

At the retained catheter count, Catheter Exchange increased HR-CTV D\textsubscript{95} by a median paired change of 2.1 cGy and reduced V\textsubscript{150} and V\textsubscript{200} by 1.0 and 1.7 percentage points, respectively, relative to Pruning (all \emph{P} < 0.01). Bladder, rectum, and colon D\textsubscript{2cc}, urethral D\textsubscript{0.1cc}, and total dwell time also decreased (all \emph{P} < 0.01), whereas small-intestine D\textsubscript{2cc} did not differ significantly. Fine-Tuning provided additional reductions in rectum D\textsubscript{2cc} (median paired change, \ensuremath{-}1.3 cGy; \emph{P} = 0.001) and the dwell-time smoothness penalty (\emph{P} < 0.001), with no significant changes in HR-CTV metrics, total dwell time, or active dwell-position count. Tables 2 and 3 present the plan metrics across optimization stages and the paired changes between consecutive stages, respectively.

\begin{table}[H]
\centering
\begin{minipage}{\linewidth}
\compactformat\let\footnotesize\compactformat\setlength{\parskip}{0pt}
\setlength{\tabcolsep}{3pt}
\renewcommand{\arraystretch}{1}
\caption{Changes between consecutive optimization stages.}\label{tab:3}
\begin{tabularx}{\linewidth}{@{}>{\raggedright\arraybackslash}p{0.14\linewidth}*{4}{>{\raggedright\arraybackslash}X}>{\raggedright\arraybackslash}p{0.085\linewidth}@{}}
\specialrule{0.5pt}{1.5pt}{1.5pt}
\textbf{\textbf{Metric}} & \textbf{\textbf{Group-Sparsity}\newline \textbf{vs.}\newline \textbf{Clinical}} & \textbf{\textbf{Pruning}\newline \textbf{vs.}\newline \textbf{Group-Sparsity}} & \textbf{\textbf{Catheter Exchange}\newline \textbf{vs.}\newline \textbf{Pruning}} & \textbf{\textbf{Fine-Tuning}\newline \textbf{vs.}\newline \textbf{Catheter Exchange}} & \textbf{\textbf{Patients/\newline fractions}} \\
\specialrule{1.5pt}{1.5pt}{1.5pt}
Number of catheters & \mbox{\ensuremath{-}1 (\ensuremath{-}3, 1)}\newline \mbox{\emph{P} = 0.016} & \mbox{\ensuremath{-}1 (\ensuremath{-}1, 0)}\newline \mbox{\emph{P} < 0.001} & \mbox{0 (0, 0)}\newline \mbox{---} & \mbox{0 (0, 0)}\newline \mbox{---} & 26/94 \\[1.5pt]
\cmidrule[0.5pt]{2-6}
HR-CTV D\textsubscript{95}, cGy & \mbox{\ensuremath{-}3.8 (\ensuremath{-}12.4, 3.6)}\newline \mbox{\emph{P} = 0.007} & \mbox{\ensuremath{-}1.7 (\ensuremath{-}7.8, 0.8)}\newline \mbox{\emph{P} = 0.001} & \mbox{2.1 (\ensuremath{-}1.0, 6.6)}\newline \mbox{\emph{P} = 0.001} & \mbox{0.5 (\ensuremath{-}2.1, 2.5)}\newline \mbox{\emph{P} = 0.770} & 26/94 \\[1.5pt]
\cmidrule[0.5pt]{2-6}
HR-CTV V\textsubscript{150}, \% & \mbox{0.0 (\ensuremath{-}1.8, 2.5)}\newline \mbox{\emph{P} = 0.144} & \mbox{1.0 (0.0, 2.4)}\newline \mbox{\emph{P} < 0.001} & \mbox{\ensuremath{-}1.0 (\ensuremath{-}2.6, 0.0)}\newline \mbox{\emph{P} = 0.003} & \mbox{0.1 (\ensuremath{-}0.4, 0.8)}\newline \mbox{\emph{P} = 0.424} & 26/94 \\[1.5pt]
\cmidrule[0.5pt]{2-6}
HR-CTV V\textsubscript{200}, \% & \mbox{0.3 (\ensuremath{-}2.8, 4.1)}\newline \mbox{\emph{P} = 0.424} & \mbox{1.6 (0.0, 3.3)}\newline \mbox{\emph{P} < 0.001} & \mbox{\ensuremath{-}1.7 (\ensuremath{-}4.2, 0.0)}\newline \mbox{\emph{P} < 0.001} & \mbox{0.0 (\ensuremath{-}0.9, 0.8)}\newline \mbox{\emph{P} = 0.424} & 26/94 \\[1.5pt]
\cmidrule[0.5pt]{2-6}
Bladder D\textsubscript{2cc}, cGy & \mbox{\ensuremath{-}0.7 (\ensuremath{-}12.4, 9.3)}\newline \mbox{\emph{P} = 0.770} & \mbox{4.3 (\ensuremath{-}0.8, 10.2)}\newline \mbox{\emph{P} = 0.002} & \mbox{\ensuremath{-}5.8 (\ensuremath{-}15.6, 0.0)}\newline \mbox{\emph{P} = 0.001} & \mbox{\ensuremath{-}0.3 (\ensuremath{-}3.6, 1.5)}\newline \mbox{\emph{P} = 0.016} & 26/94 \\[1.5pt]
\cmidrule[0.5pt]{2-6}
Rectum D\textsubscript{2cc}, cGy & \mbox{1.1 (\ensuremath{-}7.1, 10.6)}\newline \mbox{\emph{P} = 0.869} & \mbox{2.5 (\ensuremath{-}1.3, 9.9)}\newline \mbox{\emph{P} = 0.001} & \mbox{\ensuremath{-}3.1 (\ensuremath{-}10.9, 0.0)}\newline \mbox{\emph{P} = 0.002} & \mbox{\ensuremath{-}1.3 (\ensuremath{-}4.3, 0.4)}\newline \mbox{\emph{P} = 0.001} & 26/94 \\[1.5pt]
\cmidrule[0.5pt]{2-6}
Colon D\textsubscript{2cc}, cGy & \mbox{\ensuremath{-}2.9 (\ensuremath{-}11.4, 3.8)}\newline \mbox{\emph{P} = 0.064} & \mbox{1.9 (\ensuremath{-}0.1, 5.8)}\newline \mbox{\emph{P} = 0.003} & \mbox{\ensuremath{-}3.2 (\ensuremath{-}8.4, \ensuremath{-}0.5)}\newline \mbox{\emph{P} < 0.001} & \mbox{\ensuremath{-}0.1 (\ensuremath{-}1.7, 0.4)}\newline \mbox{\emph{P} = 0.053} & 19/57 \\[1.5pt]
\cmidrule[0.5pt]{2-6}
Small intestine D\textsubscript{2cc}, cGy & \mbox{\ensuremath{-}8.0 (\ensuremath{-}15.8, 1.2)}\newline \mbox{\emph{P} = 0.213} & \mbox{2.3 (0.0, 3.9)}\newline \mbox{\emph{P} = 0.003} & \mbox{\ensuremath{-}2.8 (\ensuremath{-}6.6, 0.0)}\newline \mbox{\emph{P} = 0.022} & \mbox{\ensuremath{-}0.4 (\ensuremath{-}2.2, 0.1)}\newline \mbox{\emph{P} = 0.575} & 11/19 \\[1.5pt]
\cmidrule[0.5pt]{2-6}
Urethra D\textsubscript{0.1cc}, cGy & \mbox{\ensuremath{-}10.3 (\ensuremath{-}26.1, \ensuremath{-}3.5)}\newline \mbox{\emph{P} < 0.001} & \mbox{1.4 (\ensuremath{-}0.9, 6.1)}\newline \mbox{\emph{P} = 0.002} & \mbox{\ensuremath{-}4.3 (\ensuremath{-}10.4, 0.0)}\newline \mbox{\emph{P} = 0.003} & \mbox{\ensuremath{-}0.9 (\ensuremath{-}4.4, 1.4)}\newline \mbox{\emph{P} = 0.074} & 25/90 \\[1.5pt]
\cmidrule[0.5pt]{2-6}
Total dwell time, s & \mbox{0.7 (\ensuremath{-}8.1, 14.2)}\newline \mbox{\emph{P} = 0.191} & \mbox{4.7 (0.2, 9.5)}\newline \mbox{\emph{P} < 0.001} & \mbox{\ensuremath{-}1.0 (\ensuremath{-}9.3, 2.6)}\newline \mbox{\emph{P} = 0.008} & \mbox{\ensuremath{-}0.2 (\ensuremath{-}2.7, 1.7)}\newline \mbox{\emph{P} = 0.049} & 26/94 \\[1.5pt]
\cmidrule[0.5pt]{2-6}
Number of active dwell positions & \mbox{---} & \mbox{---} & \mbox{---} & \mbox{0.0 (0.0, 1.0)}\newline \mbox{\emph{P} = 0.047} & 26/94 \\[1.5pt]
\cmidrule[0.5pt]{2-6}
Dwell-time smoothness penalty & \mbox{---} & \mbox{---} & \mbox{---} & \mbox{\ensuremath{-}2.2 (\ensuremath{-}9.0, \ensuremath{-}0.4)}\newline \mbox{\emph{P} < 0.001} & 26/94 \\[1.5pt]
\specialrule{1.5pt}{1.5pt}{1.5pt}
\end{tabularx}
\par\vspace{3pt}
\noindent Note: Comparison cells show the median paired change (first quartile, third quartile), with \emph{P} values below. Changes between automated stages were calculated at the fraction-plan level as the later-stage value minus the preceding-stage value. For Group-Sparsity versus Clinical, changes were calculated as the Group-Sparsity value minus the clinical-plan value. \emph{P} values were obtained from two-sided Wilcoxon signed-rank tests on patient-level means, with significance assessed at \emph{P} < 0.01. HR-CTV, high-risk clinical target volume.\par
\end{minipage}
\end{table}

Median total optimization time was 510.6 s. Median times for Group-Sparsity, Catheter Exchange, and Fine-Tuning were 230.1, 86.7, and 163.8 s, respectively. FISTA stopped at the active-set stability criterion in 92 plans and at the HR-CTV coverage safeguard in two; none reached 20,000 iterations. All Catheter Exchange runs stopped before 1,000 iterations. Fine-Tuning reached 1,000 iterations in one plan. Table~\ref{tab:4} summarizes the iteration counts and computation times across optimization stages.

\begin{table}[H]
\centering
\begin{minipage}{\linewidth}
\compactformat\let\footnotesize\compactformat\setlength{\parskip}{0pt}
\setlength{\tabcolsep}{3pt}
\renewcommand{\arraystretch}{1}
\caption{Optimization iterations and computation time.}\label{tab:4}
\begin{tabularx}{\linewidth}{@{}>{\raggedright\arraybackslash}p{0.26\linewidth}*{2}{>{\raggedright\arraybackslash}X}@{}}
\specialrule{0.5pt}{1.5pt}{1.5pt}
\textbf{\textbf{Metric}} & \textbf{\textbf{Iterations}} & \textbf{\textbf{Runtime, s}} \\
\specialrule{1.5pt}{1.5pt}{1.5pt}
Candidate generation & --- & 3.4 (2.5, 4.5) \\
Initial pool construction & --- & 3.0 (2.4, 4.3) \\
Group-Sparsity & 4434.5 (3523.3, 6291.3) & 230.1 (151.0, 371.1) \\
Pruning & --- & 0.9 (0.4, 1.4) \\
Catheter Exchange & 176 (129.5, 240.0) & 86.7 (64.0, 122.6) \\
Fine-Tuning & 312.5 (253.0, 444.0) & 163.8 (115.2, 230.1) \\
Total optimization & --- & 510.6 (405.9, 734.5) \\
\specialrule{1.5pt}{1.5pt}{1.5pt}
\end{tabularx}
\par\vspace{3pt}
\noindent Note: Values are presented as medians (first quartile, third quartile).\par
\end{minipage}
\end{table}

\FloatBarrier
\subsection{Ablation Studies}\label{sec:3-3}

All 94 plans completed the workflow under each evaluated configuration. The analyses examined the effects of the azimuthal constraint, initial-pool construction, and sparsity-weight configuration on the final catheter configurations and plan quality.

\subsubsection{Azimuthal Constraint}\label{sec:3-3-1}

Disabling the azimuthal constraint reduced the median initial catheter pool from 62 to 37.5 catheters (median paired change, \ensuremath{-}25; \emph{P} < 0.001), without a significant change in catheter count after Group-Sparsity or in the final plan. Final HR-CTV V\textsubscript{200} increased by 1.2 percentage points (\emph{P} = 0.003), and total optimization time decreased by 89.5 s (\emph{P} = 0.007). Further stage-specific timing analysis showed that FISTA runtime was shorter without the azimuthal constraint than in Reference (median paired change, \ensuremath{-}60.8 s; \emph{P} = 0.001). Other final dose metrics and total dwell time did not differ significantly. Table~\ref{tab:5} summarizes the ablation results for the azimuthal constraint. Four fraction plans from different patients were selected for visual comparison; each had at least six final catheters in both configurations, with the counts differing by no more than two. Figure~\ref{fig:7} shows the corresponding catheter configurations with and without the azimuthal constraint.

\begin{table}[H]
\centering
\begin{minipage}{\linewidth}
\compactformat\let\footnotesize\compactformat\setlength{\parskip}{0pt}
\setlength{\tabcolsep}{3pt}
\renewcommand{\arraystretch}{1}
\caption{Ablation results for the azimuthal constraint.}\label{tab:5}
\begin{tabularx}{\linewidth}{@{}>{\raggedright\arraybackslash}p{0.18\linewidth}*{3}{>{\raggedright\arraybackslash}X}>{\raggedright\arraybackslash}p{0.07\linewidth}>{\raggedright\arraybackslash}p{0.10\linewidth}@{}}
\specialrule{0.5pt}{1.5pt}{1.5pt}
\textbf{\textbf{Metric}} & \textbf{\textbf{Reference}} & \textbf{\textbf{No azimuthal constraint}} & \textbf{\textbf{Paired change}} & \textbf{\textbf{\emph{P} value}} & \textbf{\textbf{Patients/\newline fractions}} \\
\specialrule{1.5pt}{1.5pt}{1.5pt}
Initial catheter pool size & 62 (56, 68) & 37.5 (34, 40) & \ensuremath{-}25 (\ensuremath{-}28, \ensuremath{-}21) & <0.001 & 26/94 \\
Number of active catheters after Group-Sparsity & 5 (4, 6) & 5 (3, 7) & 0 (\ensuremath{-}1, 1) & 0.538 & 26/94 \\
Final catheter number & 4 (3, 5) & 4 (3, 5) & 0 (\ensuremath{-}1, 1) & 0.437 & 26/94 \\
HR-CTV D\textsubscript{95}, cGy & 538.2 (527.9, 543.8) & 537.9 (530.0, 545.8) & 0.9 (\ensuremath{-}4.3, 6.7) & 0.603 & 26/94 \\
HR-CTV V\textsubscript{150}, \% & 55.9 (53.2, 58.0) & 56.8 (54.1, 58.4) & 0.6 (\ensuremath{-}0.5, 2.2) & 0.023 & 26/94 \\
HR-CTV V\textsubscript{200}, \% & 31.0 (25.2, 35.0) & 33.2 (28.3, 36.7) & 1.2 (\ensuremath{-}0.5, 4.5) & 0.003 & 26/94 \\
Bladder D\textsubscript{2cc}, cGy & 423.9 (325.4, 463.4) & 422.7 (323.8, 463.0) & \ensuremath{-}0.6 (\ensuremath{-}5.8, 4.6) & 0.469 & 26/94 \\
Rectum D\textsubscript{2cc}, cGy & 363.3 (325.6, 393.3) & 359.8 (324.0, 397.9) & \ensuremath{-}0.8 (\ensuremath{-}4.7, 4.9) & 0.409 & 26/94 \\
Colon D\textsubscript{2cc}, cGy & 183.4 (101.0, 314.9) & 182.5 (102.1, 309.0) & 0.2 (\ensuremath{-}3.0, 2.6) & 0.748 & 19/57 \\
Small intestine D\textsubscript{2cc}, cGy & 121.0 (102.3, 183.4) & 118.0 (101.4, 186.7) & 1.2 (\ensuremath{-}1.3, 3.4) & 0.155 & 11/19 \\
Urethra D\textsubscript{0.1cc}, cGy & 314.0 (186.3, 420.3) & 314.3 (193.2, 421.3) & 1.1 (\ensuremath{-}4.0, 8.2) & 0.397 & 25/90 \\
Total dwell time, s & 244.0 (177.2, 301.3) & 239.9 (181.6, 313.0) & 4.0 (\ensuremath{-}6.4, 10.5) & 0.144 & 26/94 \\
Total optimization time, s & 510.6 (405.9, 734.5) & 443.7 (372.8, 584.6) & \ensuremath{-}89.5 (\ensuremath{-}190.4, 37.1) & 0.007 & 26/94 \\
\specialrule{1.5pt}{1.5pt}{1.5pt}
\end{tabularx}
\par\vspace{3pt}
\noindent Note: Values are presented as medians (first quartile, third quartile), except for \emph{P} values and patient/fraction counts. Paired changes were calculated at the fraction-plan level by subtracting the Reference value from the value obtained without the azimuthal constraint. \emph{P} values were obtained from two-sided Wilcoxon signed-rank tests on patient-level means, with significance assessed at \emph{P} < 0.01. HR-CTV, high-risk clinical target volume.\par
\end{minipage}
\end{table}

\FloatBarrier
\begin{figure}[H]
\centering
\includegraphics[width=0.92\linewidth,height=0.48\textheight,keepaspectratio]{figures20261114/figure7.jpg}
\caption{Final catheter configurations with and without the azimuthal constraint. Four selected fraction plans are shown in side and axial views of the template. HR-CTV volumes and final catheter counts are indicated for each plan. White arrows indicate projected crossings between catheter trajectories in the configurations without the azimuthal constraint.}\label{fig:7}
\end{figure}

\subsubsection{Initial Pool and Sparsity Weights}\label{sec:3-3-2}

With identical initial catheter pools, Uniform-weight yielded a higher median final catheter count than Reference (five versus four; median paired change, +1; \emph{P} = 0.002). This higher count was accompanied by lower HR-CTV V\textsubscript{200}, colon D\textsubscript{2cc}, and total dwell time, with median paired reductions of 1.6 percentage points, 1.8 cGy, and 5.3 s, respectively (all \emph{P} < 0.01). The difference in active catheter count immediately after Group-Sparsity did not meet the prespecified significance threshold (\emph{P} = 0.013).

\begin{table}[H]
\centering
\begin{minipage}{\linewidth}
\compactformat\let\footnotesize\compactformat\setlength{\parskip}{0pt}
\setlength{\tabcolsep}{3pt}
\renewcommand{\arraystretch}{1}
\caption{Plan metrics by initial-pool construction and sparsity-weight configuration.}\label{tab:6}
\begin{tabularx}{\linewidth}{@{}>{\raggedright\arraybackslash}p{0.26\linewidth}*{3}{>{\raggedright\arraybackslash}X}@{}}
\specialrule{0.5pt}{1.5pt}{1.5pt}
\textbf{\textbf{Metric}} & \textbf{\textbf{Reference}} & \textbf{\textbf{Uniform-weight}} & \textbf{\textbf{Random-pool}} \\
\specialrule{1.5pt}{1.5pt}{1.5pt}
Initial catheter pool size & 62 (56, 68) & 62 (56, 68) & 63 (58, 66) \\
Number of active catheters after Group-Sparsity & 5 (4, 6) & 6 (5, 7) & 5 (4, 7) \\
Final catheter number & 4 (3, 5) & 5 (4, 6) & 5 (4, 6) \\
HR-CTV D\textsubscript{95}, cGy & 538.2 (527.9, 543.8) & 537.0 (525.7, 545.5) & 539.5 (527.7, 545.8) \\
HR-CTV V\textsubscript{150}, \% & 55.9 (53.2, 58.0) & 55.4 (52.6, 57.4) & 55.4 (52.1, 57.7) \\
HR-CTV V\textsubscript{200}, \% & 31.0 (25.2, 35.0) & 28.7 (23.9, 32.6) & 28.6 (24.9, 33.8) \\
Bladder D\textsubscript{2cc}, cGy & 423.9 (325.4, 463.4) & 418.6 (320.4, 459.9) & 417.3 (317.9, 462.5) \\
Rectum D\textsubscript{2cc}, cGy & 363.3 (325.6, 393.3) & 358.2 (319.2, 393.4) & 359.2 (321.8, 392.5) \\
Colon D\textsubscript{2cc}, cGy & 183.4 (101.0, 314.9) & 183.6 (99.1, 305.5) & 183.1 (98.7, 309.2) \\
Small intestine D\textsubscript{2cc}, cGy & 121.0 (102.3, 183.4) & 117.3 (101.5, 187.5) & 120.4 (100.8, 179.3) \\
Urethra D\textsubscript{0.1cc}, cGy & 314.0 (186.3, 420.3) & 314.2 (186.1, 419.2) & 308.3 (186.2, 417.3) \\
Total dwell time, s & 244.0 (177.2, 301.3) & 237.4 (172.6, 301.5) & 230.9 (174.1, 301.0) \\
Total optimization time, s & 510.6 (405.9, 734.5) & 555.9 (436.3, 714.3) & 508.6 (412.6, 703.9) \\
\specialrule{1.5pt}{1.5pt}{1.5pt}
\end{tabularx}
\par\vspace{3pt}
\noindent Note: Values are presented as medians (first quartile, third quartile). HR-CTV, high-risk clinical target volume.\par
\end{minipage}
\end{table}

Uniform-weight and Random-pool did not differ significantly in catheter count, dosimetry, total dwell time, or optimization time. Relative to Reference, Random-pool increased the median final catheter count from four to five (\emph{P} = 0.009) and decreased total dwell time by 3.9 s (\emph{P} = 0.005), without significant changes in final dose metrics or optimization time. Tables 6 and 7 present the plan metrics for the three configurations and their paired comparisons, respectively.

\begin{table}[H]
\centering
\begin{minipage}{\linewidth}
\compactformat\let\footnotesize\compactformat\setlength{\parskip}{0pt}
\setlength{\tabcolsep}{3pt}
\renewcommand{\arraystretch}{1}
\caption{Paired comparisons of initial-pool construction and sparsity-weight configurations.}\label{tab:7}
\begin{tabularx}{\linewidth}{@{}>{\raggedright\arraybackslash}p{0.20\linewidth}*{3}{>{\raggedright\arraybackslash}X}>{\raggedright\arraybackslash}p{0.09\linewidth}@{}}
\specialrule{0.5pt}{1.5pt}{1.5pt}
\textbf{\textbf{Metric}} & \textbf{\textbf{Uniform-weight}\newline \textbf{vs.}\newline \textbf{Reference}} & \textbf{\textbf{Random-pool}\newline \textbf{vs.}\newline \textbf{Uniform-weight}} & \textbf{\textbf{Random-pool}\newline \textbf{vs.}\newline \textbf{Reference}} & \textbf{\textbf{Patients/\newline fractions}} \\
\specialrule{1.5pt}{1.5pt}{1.5pt}
Initial catheter pool size & \mbox{0 (0, 0)}\newline \mbox{---} & \mbox{0 (\ensuremath{-}2, 3)}\newline \mbox{\emph{P} = 0.516} & \mbox{0 (\ensuremath{-}2, 3)}\newline \mbox{\emph{P} = 0.516} & 26/94 \\[1pt]
\cmidrule[0.5pt]{2-5}
Number of active catheters after Group-Sparsity & \mbox{1 (0, 2)}\newline \mbox{\emph{P} = 0.013} & \mbox{0 (\ensuremath{-}2, 1)}\newline \mbox{\emph{P} = 0.201} & \mbox{0 (\ensuremath{-}1, 2)}\newline \mbox{\emph{P} = 0.082} & 26/94 \\[1pt]
\cmidrule[0.5pt]{2-5}
Final catheter number & \mbox{1 (0, 2)}\newline \mbox{\emph{P} = 0.002} & \mbox{0 (\ensuremath{-}2, 1)}\newline \mbox{\emph{P} = 0.269} & \mbox{1 (0, 2)}\newline \mbox{\emph{P} = 0.009} & 26/94 \\[1pt]
\cmidrule[0.5pt]{2-5}
HR-CTV D\textsubscript{95}, cGy & \mbox{\ensuremath{-}1.5 (\ensuremath{-}7.4, 5.6)}\newline \mbox{\emph{P} = 0.809} & \mbox{1.0 (\ensuremath{-}4.6, 8.3)}\newline \mbox{\emph{P} = 0.107} & \mbox{0.4 (\ensuremath{-}5.0, 6.1)}\newline \mbox{\emph{P} = 0.849} & 26/94 \\[1pt]
\cmidrule[0.5pt]{2-5}
HR-CTV V\textsubscript{150}, \% & \mbox{0.0 (\ensuremath{-}2.3, 1.8)}\newline \mbox{\emph{P} = 0.469} & \mbox{0.4 (\ensuremath{-}1.7, 2.2)}\newline \mbox{\emph{P} = 0.137} & \mbox{\ensuremath{-}0.3 (\ensuremath{-}1.9, 2.0)}\newline \mbox{\emph{P} = 0.732} & 26/94 \\[1pt]
\cmidrule[0.5pt]{2-5}
HR-CTV V\textsubscript{200}, \% & \mbox{\ensuremath{-}1.6 (\ensuremath{-}4.6, 1.3)}\newline \mbox{\emph{P} = 0.004} & \mbox{0.7 (\ensuremath{-}3.3, 4.3)}\newline \mbox{\emph{P} = 0.182} & \mbox{\ensuremath{-}1.1 (\ensuremath{-}5.9, 1.9)}\newline \mbox{\emph{P} = 0.030} & 26/94 \\[1pt]
\cmidrule[0.5pt]{2-5}
Bladder D\textsubscript{2cc}, cGy & \mbox{\ensuremath{-}1.8 (\ensuremath{-}7.4, 2.0)}\newline \mbox{\emph{P} = 0.019} & \mbox{\ensuremath{-}0.1 (\ensuremath{-}5.7, 5.1)}\newline \mbox{\emph{P} = 0.439} & \mbox{\ensuremath{-}1.1 (\ensuremath{-}7.6, 3.7)}\newline \mbox{\emph{P} = 0.269} & 26/94 \\[1pt]
\cmidrule[0.5pt]{2-5}
Rectum D\textsubscript{2cc}, cGy & \mbox{\ensuremath{-}1.7 (\ensuremath{-}5.4, 2.1)}\newline \mbox{\emph{P} = 0.131} & \mbox{\ensuremath{-}1.1 (\ensuremath{-}5.7, 5.0)}\newline \mbox{\emph{P} = 0.439} & \mbox{\ensuremath{-}1.4 (\ensuremath{-}7.0, 3.9)}\newline \mbox{\emph{P} = 0.209} & 26/94 \\[1pt]
\cmidrule[0.5pt]{2-5}
Colon D\textsubscript{2cc}, cGy & \mbox{\ensuremath{-}1.8 (\ensuremath{-}4.8, 0.1)}\newline \mbox{\emph{P} = 0.008} & \mbox{0.0 (\ensuremath{-}2.0, 3.5)}\newline \mbox{\emph{P} = 0.027} & \mbox{\ensuremath{-}1.4 (\ensuremath{-}4.6, 0.9)}\newline \mbox{\emph{P} = 0.064} & 19/57 \\[1pt]
\cmidrule[0.5pt]{2-5}
Small intestine D\textsubscript{2cc}, cGy & \mbox{\ensuremath{-}0.3 (\ensuremath{-}3.5, 2.0)}\newline \mbox{\emph{P} = 0.959} & \mbox{\ensuremath{-}0.7 (\ensuremath{-}2.0, 1.9)}\newline \mbox{\emph{P} = 0.859} & \mbox{\ensuremath{-}0.8 (\ensuremath{-}2.9, 2.3)}\newline \mbox{\emph{P} = 0.534} & 11/19 \\[1pt]
\cmidrule[0.5pt]{2-5}
Urethra D\textsubscript{0.1cc}, cGy & \mbox{\ensuremath{-}2.1 (\ensuremath{-}6.5, 3.3)}\newline \mbox{\emph{P} = 0.668} & \mbox{\ensuremath{-}0.5 (\ensuremath{-}6.9, 4.7)}\newline \mbox{\emph{P} = 0.339} & \mbox{\ensuremath{-}3.3 (\ensuremath{-}7.6, 4.3)}\newline \mbox{\emph{P} = 0.150} & 25/90 \\[1pt]
\cmidrule[0.5pt]{2-5}
Total dwell time, s & \mbox{\ensuremath{-}5.3 (\ensuremath{-}16.0, 2.0)}\newline \mbox{\emph{P} < 0.001} & \mbox{\ensuremath{-}0.4 (\ensuremath{-}7.8, 7.6)}\newline \mbox{\emph{P} = 0.551} & \mbox{\ensuremath{-}3.9 (\ensuremath{-}16.7, 3.9)}\newline \mbox{\emph{P} = 0.005} & 26/94 \\[1pt]
\cmidrule[0.5pt]{2-5}
Total optimization time, s & \mbox{26.3 (\ensuremath{-}123.6, 136.2)}\newline \mbox{\emph{P} = 0.675} & \mbox{\ensuremath{-}20.3 (\ensuremath{-}176.5, 113.6)}\newline \mbox{\emph{P} = 0.849} & \mbox{\ensuremath{-}0.6 (\ensuremath{-}107.0, 86.9)}\newline \mbox{\emph{P} = 0.770} & 26/94 \\[1pt]
\specialrule{1.5pt}{1.5pt}{1.5pt}
\end{tabularx}
\par\vspace{3pt}
\noindent Note: Comparison cells show the median paired change (first quartile, third quartile), with \emph{P} values below. Changes were calculated at the fraction-plan level as the first-listed configuration minus the second-listed configuration in each column heading. \emph{P} values were obtained from two-sided Wilcoxon signed-rank tests on patient-level means, with significance assessed at \emph{P} < 0.01. HR-CTV, high-risk clinical target volume.\par
\end{minipage}
\end{table}

\section{Discussion}\label{sec:4}

The proposed workflow reduced the median catheter number from 6 in the clinical plans to 4 in the final automated plans; 78 of 94 fraction plans used fewer catheters. After normalization to the same HR-CTV D\textsubscript{90}, the automated plans showed a small but significant decrease in D\textsubscript{95}, no significant differences in V\textsubscript{150} or V\textsubscript{200}, and lower colon D\textsubscript{2cc} and urethral D\textsubscript{0.1cc}. These findings suggest that reduced catheter use was accompanied by a small trade-off in HR-CTV D\textsubscript{95} and dose reductions in selected OARs.

Cases with increased catheter counts or particularly large reductions warrant further discussion. Five automated plans used one or two additional catheters, possibly reflecting restrictions on accessible trajectories imposed by the predefined sampling regions and geometric constraints. Conversely, the four largest reductions occurred in large targets (HR-CTV volumes, 105.5--246.8 cm\textsuperscript{3}), with catheter counts reduced from 9--11 in the clinical plans to 3--4 in the automated plans. In these four plans, HR-CTV D\textsubscript{95} decreased by 4.4--9.4 cGy after normalization, while 12 of 17 available OAR dose comparisons showed reductions. Among these four plans, the two largest OAR dose increases occurred in the same fraction, with colon D\textsubscript{2cc} rising from 382.8 to 399.6 cGy and small-intestine D\textsubscript{2cc} from 268.6 to 305.8 cGy. These findings support the dosimetric feasibility of substantial catheter reduction for these large targets within the planning model. However, when implanted catheters deviate from their planned trajectories because of needle bending or placement errors,\cite{ref13} having fewer catheters may limit the ability to compensate through dwell-time reoptimization. In a prostate HDR brachytherapy planning study, plans with more catheters were more likely to remain dosimetrically acceptable under simulated catheter displacement.\cite{ref68} Reserve guide channels could therefore be incorporated into the individualized template before implantation to provide alternative catheter trajectories if needed, as previously described for 3D-PIT.\cite{ref21}

The stage-wise analysis distinguished catheter reduction from subsequent configuration optimization. The Group-Sparsity stage yielded an initial reduction in catheter use with accompanying dosimetric trade-offs, although the catheter-number difference from the clinical plans did not meet the prespecified significance threshold. Pruning further reduced the median catheter number from five to four by removing catheters with low total dwell times, thereby decreasing the number of available dwell positions and limiting dose-shaping flexibility. Increasing dwell times at the remaining fixed positions to compensate for lost target dose may also increase dose to other target regions and nearby OARs. The increases in target high-dose volumes and OAR doses observed after Pruning are consistent with this effect. With the catheter count fixed, Catheter Exchange can replace trajectories to provide dwell positions that deliver dose more effectively to previously undercovered target regions. After dwell-time reoptimization, these regions can receive dose from the new positions, reducing the need for compensatory increases in dwell times elsewhere and helping improve D\textsubscript{95} while reducing target high-dose volumes and OAR doses. Fine-Tuning adjusted catheter endpoints to explore trajectories beyond the sampled candidate set, but its additional effect on conventional dose--volume indices after Catheter Exchange was limited. HR-CTV metrics did not differ significantly, and rectum D\textsubscript{2cc} was the only OAR metric with a significant additional reduction. The dwell-time smoothness penalty decreased without a significant change in active dwell-position count, indicating smoother optimized dwell-time distributions. Whether these smoother distributions improve robustness to source-positioning uncertainties requires further evaluation. In this retrospective analysis, the median computation time was 337.2 s (5.6 min) through Catheter Exchange and 510.6 s (8.5 min) for the full workflow. For future pre-implantation planning, the improvement in dwell-time smoothness with Fine-Tuning should be weighed against the additional computation time.

The ablation results suggest that differential sparsity weighting can help reduce the final catheter count when both configurations use the same dose-guided initial pool. Dose-guided core selection provided the basis for differential sparsity weighting in Reference by distinguishing core catheters from randomly supplemented catheters. Relative to Uniform-weight, Reference assigned a lower sparsity multiplier to the core group and a higher multiplier to the supplementary group to favor core catheter retention. The reduction in catheter count came with trade-offs in dosimetry and total dwell time. Compared with Uniform-weight, Reference yielded a lower median final catheter count (four vs. five), with paired comparisons showing higher HR-CTV V\textsubscript{200}, higher colon D\textsubscript{2cc}, and longer total dwell time. Under uniform weighting, the dose-guided and random initial pools did not differ significantly in active catheter count after Group-Sparsity, final catheter count, dosimetry, total dwell time, or optimization time.

The azimuthal constraint limits entry--target azimuthal differences to reduce oblique trajectories spanning different azimuthal sectors. In the four selected plans shown in Figure~\ref{fig:7}, enabling the constraint produced more consistently aligned catheter trajectories with fewer apparent crossings in the projected views. Beyond these differences in catheter distribution, the ablation comparisons also showed significant differences in initial-pool size, total optimization time, and HR-CTV V\textsubscript{200}. The smaller initial pool obtained without the azimuthal constraint may reflect more candidate and selected trajectories following oblique paths across different azimuthal sectors, making it harder for subsequent candidates to maintain the required 2-mm separation from previously selected catheters along their entire paths. The significantly longer total optimization time in Reference may be related to its larger initial pool, which requires simultaneous optimization of dwell times for more catheters. Further stage-specific timing analysis showed a median paired change of +60.8 s in FISTA runtime for Reference relative to the configuration without the azimuthal constraint (\emph{P} = 0.001). Final catheter counts did not differ significantly between configurations, but Reference showed a median paired change of \ensuremath{-}1.2 percentage points in HR-CTV V\textsubscript{200} relative to the configuration without the azimuthal constraint (\emph{P} = 0.003), suggesting that the constraint may also help control target high-dose volumes.

This study has several limitations. It was a single-center retrospective computational replanning study using post-implantation CT images. The generalizability of the findings to broader clinical populations remains uncertain. The automated trajectories were not used to fabricate templates or guide catheter insertion, and their prospective implementation remains to be evaluated. Not all OARs were contoured in every fraction, particularly the small intestine, which was delineated in only 19 fractions from 11 patients; the corresponding results should therefore be interpreted cautiously. Because the optimization bounds were derived from each fraction's original, unrescaled clinical plan, planning with independently specified dose objectives was not evaluated. Several algorithmic parameters, including the 30\textdegree{} azimuthal threshold, the HR-CTV D\textsubscript{90} feasibility threshold set at 95\% of the corresponding original, unrescaled clinical-plan D\textsubscript{90}, the 5\% cumulative dwell-time limit for Pruning, and the perturbation magnitude used in Fine-Tuning, were selected based on clinical experience. Further sensitivity analyses are needed to assess the impact of these parameter choices on planning performance. The nonconvex sparsity formulation and heuristic stopping criteria do not guarantee a globally optimal catheter configuration. Finally, the study evaluated catheter counts, dosimetric metrics, and computation time but not procedural burden, toxicity, or clinical outcomes.

\FloatBarrier
\begin{samepage}
\section{Conclusions}\label{sec:5}

This study presents an integrated automated planning approach that combines catheter generation and initial-pool construction with adaptive catheter number determination and subsequent fixed-count spatial configuration and dwell-time optimization for individualized template-guided HDR-ISBT. In retrospective computational replanning, the framework reduced the median catheter number from six to four. After normalization to the same HR-CTV D\textsubscript{90}, the automated plans showed a small but statistically significant reduction in HR-CTV D\textsubscript{95}, no significant differences in V\textsubscript{150} or V\textsubscript{200}, and lower colon D\textsubscript{2cc} and urethral D\textsubscript{0.1cc} than the clinical plans. Stage-wise analysis showed that Pruning enabled further catheter reduction but adversely affected several dosimetric metrics; subsequent Catheter Exchange mitigated several of these changes at the retained catheter count. The framework therefore provides an automated approach for balancing catheter use and dosimetric performance, although prospective pre-implantation planning and physical validation of template fabrication and catheter insertion are required before clinical implementation.

\end{samepage}

\section*{Acknowledgments}

This work was supported by the Beijing Natural Science Foundation (No. L232037, L222104 and L222034), the Key Specialty program of Natural Science Foundation of Beijing Municipality (Z20008 to P. Jiang) and the National Natural Science Foundation of China (No. 12375359).

\section*{Conflict of Interest Statement}

The authors declare that they have no conflicts of interest in this work.

\section*{Data Availability Statement}

The data that support the findings of this study are available on request from the corresponding author.

\section*{Declaration of Generative AI}

The authors used ChatGPT for language polishing and initial content structuring. The authors critically evaluated all AI-generated output, revised it where necessary, and take full responsibility for the final manuscript.

\begingroup

\setstretch{1}

\begin{thebibliography}{99}

\bibitem{ref1} Mahantshetty U, Kalyani N, Engineer R, et al. Reirradiation using high-dose-rate brachytherapy in recurrent carcinoma of uterine cervix. Brachytherapy. 2014;13(6):548-553.

\bibitem{ref2} Ren XJ, Fu YL, Liu ZS, et al. Image-guided interstitial brachytherapy for recurrent cervical cancer after radiotherapy: A single institution experience. FRONTIERS IN ONCOLOGY. 2022;12.

\bibitem{ref3} Isohashi F, Yoshida K, Murakami N, et al. Reirradiation for recurrent gynecologic cancer using high-dose-rate brachytherapy in Japan: A multicenter survey on practice patterns and outcomes. Radiother Oncol. 2024;195.

\bibitem{ref4} Sturdza A, Viswanathan AN, Erickson B, et al. American Brachytherapy Society working group report on the patterns of care and a literature review of reirradiation for gynecologic cancers. Brachytherapy. 2020;19(2):127-138.

\bibitem{ref5} Fokdal L, Sturdza A, Mazeron R, et al. Image guided adaptive brachytherapy with combined intracavitary and interstitial technique improves the therapeutic ratio in locally advanced cervical cancer: Analysis from the retro EMBRACE study. Radiother Oncol. 2016;120(3):434-440.

\bibitem{ref6} Murakami N, Ando K, Murata M, et al. An Asian multi-national multi-institutional retrospective study comparing intracavitary versus the hybrid of intracavitary and interstitial brachytherapy for locally advanced uterine cervical carcinoma. JOURNAL OF RADIATION RESEARCH. 2022;63(3):412-427.

\bibitem{ref7} Serban M, Kirisits C, de Leeuw A, et al. Ring Versus Ovoids and Intracavitary Versus Intracavitary-Interstitial Applicators in Cervical Cancer Brachytherapy: Results From the EMBRACE I Study. Int J Radiat Oncol Biol Phys. 2020;106(5):1052-1062.

\bibitem{ref8} Serban M, Fokdal L, Nielsen SK, et al. Characterization of combined intracavitary/interstitial brachytherapy including oblique needles in locally advanced cervix cancer. Brachytherapy. 2021;20(4):796-806.

\bibitem{ref9} Fokdal L, Tanderup K, Hokland SB, et al. Clinical feasibility of combined intracavitary/interstitial brachytherapy in locally advanced cervical cancer employing MRI with a tandem/ring applicator in situ and virtual preplanning of the interstitial component. Radiother Oncol. 2013;107(1):63-68.

\bibitem{ref10} Petric P, Hudej R, Hanuna O, et al. MRI-assisted cervix cancer brachytherapy pre-planning, based on application in paracervical anaesthesia: final report. Radiol Oncol. 2014;48(3):293-300.

\bibitem{ref11} Rodriguez S, Otal A, Richart J, Perez-Colatayud J, Santos M. Pre-plan technique feasibility in multi-interstitial/endocavitary perineal gynecological brachytherapy. J Contemp Brachytherapy. 2017;9(5):472-476.

\bibitem{ref12} Takagawa Y, Tsuchiya K, Izumi S, et al. A novel intra pre-plan method of interstitial needle optimization for cervical cancer patients treated with Venezia applicator brachytherapy: a technical note. J Contemp Brachytherapy. 2022;14(3):260-267.

\bibitem{ref13} Karius A, Strnad V, Lotter M, et al. Assessment of needle bending and tracking requirements for optimized needle placement in combined intracavitary/interstitial gynecologic brachytherapy. STRAHLENTHERAPIE UND ONKOLOGIE. 2026;202(1):40-51.

\bibitem{ref14} Lindegaard JC, Madsen ML, Traberg A, et al. Individualised 3D printed vaginal template for MRI guided brachytherapy in locally advanced cervical cancer. Radiother Oncol. 2016;118(1):173-175.

\bibitem{ref15} Zhao ZP, Tang XD, Mao Z, Zhao HF. The design of an individualized cylindrical vaginal applicator with oblique guide holes using 3D modeling and printing technologies. J Contemp Brachytherapy. 2019;11(5):479-487.

\bibitem{ref16} Wiebe E, Easton H, Thomas G, Barbera L, D'Alimonte L, Ravi A. Customized vaginal vault brachytherapy with computed tomography imaging-derived applicator prototyping. Brachytherapy. 2015;14(3):380-384.

\bibitem{ref17} Logar HBZ, Hudej R, Segedin B. Development and assessment of 3D-printed individual applicators in gynecological MRI-guided brachytherapy. J Contemp Brachytherapy. 2019;11(2):128-136.

\bibitem{ref18} Kudla M, Bachand F, Moore J, Batchelar D. Patient-specific cylinder templates for hybrid interstitial vaginal brachytherapy: Feasibility of automated 3-D design, 3D printing, and dosimetric outlook. Brachytherapy. 2023;22(4):468-476.

\bibitem{ref19} Sethi R, Cunha A, Mellis K, et al. Clinical applications of custom-made vaginal cylinders constructed using three-dimensional printing technology. J Contemp Brachytherapy. 2016;8(3):210-216.

\bibitem{ref20} Jiang P, Wang J, Qu A, Sun H, Zhang X, Li X. Three Dimensional Printed Individual Template Based 192ir High-Dose Rate Interstitial Brachytherapy for Central Recurrent Gynecologic Cancer: A Phase II Clinical Trial. Int J Radiat Oncol Biol Phys. 2019;105(1):E344-E344.

\bibitem{ref21} Jiang P, Qu A, Wei SH, et al. The Preliminary Results of 3-Dimensional Printed Individual Template Assisted 192Ir High-Dose Rate Interstitial Brachytherapy for Central Recurrent Gynecologic Cancer. Technol Cancer Res Treat. 2020;19.

\bibitem{ref22} Qin X, Zhang FQ, Hou XR, et al. Efficacy and safety of a 3D-printed applicator for vaginal brachytherapy in patients with central pelvic-recurrent cervical cancer after primary hysterectomy. Brachytherapy. 2022;21(2):193-201.

\bibitem{ref23} Cobussen A, Petric P, Wulff CN, et al. Clinical outcomes using a 3D printed tandem-needle-template and the EMBRACE-II planning aims for image guided adaptive brachytherapy in locally advanced cervical cancer. ACTA ONCOLOGICA. 2023;62(11):1470-1478.

\bibitem{ref24} Wang KY, Qu A, Deng XW, et al. Efficacy and safety of 3-dimensional printing noncoplanar template (3D-PNCT)-assisted high-dose-rate interstitial brachytherapy (HDR-ISBT) for reirradiation of recurrent cervical cancer: a prospective cohort. JOURNAL OF GYNECOLOGIC ONCOLOGY. 2025;36(2).

\bibitem{ref25} Yao QL, Wang KY, Liu B, et al. Integrated optimization of needle paths and dwell time for individualized template-guided interstitial brachytherapy. MEDICAL PHYSICS. 2026;53(4).

\bibitem{ref26} Shen JL, Sun SX, Xu ZW, et al. A rapid semi-automated trajectory planning method for non-coplanar interstitial brachytherapy in cervical cancer using a greedy algorithm. RADIATION ONCOLOGY. 2026;21(1).

\bibitem{ref27} Lessard E, Pouliot J. Inverse planning anatomy-based dose optimization for HDR-brachytherapy of the prostate using fast simulated annealing algorithm and dedicated objective function. MEDICAL PHYSICS. 2001;28(5):773-779.

\bibitem{ref28} Hsu ICJ, Lessard E, Weinberg V, Pouliot J. Comparison of inverse planning simulated annealing and geometrical optimization for prostate high-dose-rate brachytherapy. Brachytherapy. 2004;3(3):147-152.

\bibitem{ref29} Pokharel S, Rana S, Blikenstaff J, Sadeghi A, Prestidge B. Evaluation of hybrid inverse planning and optimization (HIPO) algorithm for optimization in real-time, high-dose-rate (HDR) brachytherapy for prostate. J Appl Clin Med Phys. 2013;14(4):96-107.

\bibitem{ref30} Lahanas M, Baltas D, Giannouli S. Global convergence analysis of fast multiobjective gradient-based dose optimization algorithms for high-dose-rate brachytherapy. Phys Med Biol. 2003;48(5):599-617.

\bibitem{ref31} Yao R, Templeton AK, Liao YX, Turian JV, Kiel KD, Chu JCH. Optimization for high-dose-rate brachytherapy of cervical cancer with adaptive simulated annealing and gradient descent. Brachytherapy. 2014;13(4):352-360.

\bibitem{ref32} Siauw T, Cunha A, Atamtürk A, Hsu IC, Pouliot J, Goldberg K. IPIP: A new approach to inverse planning for HDR brachytherapy by directly optimizing dosimetric indices. MEDICAL PHYSICS. 2011;38(7):4045-4051.

\bibitem{ref33} Gorissen BL, den Hertog D, Hoffmann AL. Mixed integer programming improves comprehensibility and plan quality in inverse optimization of prostate HDR brachytherapy. Phys Med Biol. 2013;58(4):1041-1057.

\bibitem{ref34} Deist TM, Gorissen BL. High-dose-rate prostate brachytherapy inverse planning on dose-volume criteria by simulated annealing. Phys Med Biol. 2016;61(3):1155-1170.

\bibitem{ref35} Morén B, Larsson T, Tedgren Å. Mathematical optimization of high dose-rate brachytherapy-derivation of a linear penalty model from a dose-volume model. Phys Med Biol. 2018;63(6).

\bibitem{ref36} Morén B, Larsson T, Tedgren Å. An extended dose-volume model in high dose-rate brachytherapy - Using mean-tail-dose to reduce tumor underdosage. MEDICAL PHYSICS. 2019;46(6):2556-2566.

\bibitem{ref37} Guthier CV, Damato AL, Viswanathan AN, Hesser JW, Cormack RA. A fast multitarget inverse treatment planning strategy optimizing dosimetric measures for high-dose-rate (HDR) brachytherapy. MEDICAL PHYSICS. 2017;44(9):4452-4462.

\bibitem{ref38} Cui SY, Després P, Beaulieu L. A multi-criteria optimization approach for HDR prostate brachytherapy: I. Pareto surface approximation. Phys Med Biol. 2018;63(20).

\bibitem{ref39} Cui SY, Després P, Beaulieu L. A multi-criteria optimization approach for HDR prostate brachytherapy: II. Benchmark against clinical plans. Phys Med Biol. 2018;63(20).

\bibitem{ref40} Bélanger C, Cui SY, Ma YZ, Després P, Cunha JAM, Beaulieu L. A GPU-based multi-criteria optimization algorithm for HDR brachytherapy. Phys Med Biol. 2019;64(10).

\bibitem{ref41} Breedveld S, Bennan ABA, Aluwini S, Schaart DR, Kolkman-Deurloo IKK, Heijmen BM. Fast automated multi-criteria planning for HDR brachytherapy explored for prostate cancer. Phys Med Biol. 2019;64(20).

\bibitem{ref42} Luong NH, Alderliesten T, Pieters BR, Bel A, Niatsetski Y, Bosman PAN. Fast and insightful bi-objective optimization for prostate cancer treatment planning with high-dose-rate brachytherapy. Appl Soft Comput. 2019;84.

\bibitem{ref43} Oud M, Kolkman-Deurloo IK, Mens JW, et al. Fast and fully-automated multi-criterial treatment planning for adaptive HDR brachytherapy for locally advanced cervical cancer. Radiother Oncol. 2020;148:143-150.

\bibitem{ref44} Cunha A, Siauw T, Hsu IC, Pouliot J. A method for restricting intracatheter dwell time variance in high-dose-rate brachytherapy plan optimization. Brachytherapy. 2016;15(2):246-251.

\bibitem{ref45} Shen CY, Gonzalez Y, Klages P, et al. Intelligent inverse treatment planning via deep reinforcement learning, a proof-of-principle study in high dose-rate brachytherapy for cervical cancer. Phys Med Biol. 2019;64(11).

\bibitem{ref46} Mitra A, Moore LC, Kallis K, et al. Rapid Artificial Intelligence Autoplanning Rivals Manual Expert Planning for Cervical Brachytherapy. Practical Radiation Oncology. 2026.

\bibitem{ref47} Siauw T, Cunha A, Berenson D, et al. NPIP: A skew line needle configuration optimization system for HDR brachytherapy. MEDICAL PHYSICS. 2012;39(7):4339-4346.

\bibitem{ref48} Poulin E, Fekete CAC, Létourneau M, Fenster A, Pouliot J, Beaulieu L. Adaptation of the CVT algorithm for catheter optimization in high dose rate brachytherapy. MEDICAL PHYSICS. 2013;40(11).

\bibitem{ref49} Poulin E, Varfalvy N, Aubin S, Beaulieu L. Comparison of dose and catheter optimization algorithms in prostate high-dose-rate brachytherapy. Brachytherapy. 2016;15(1):102-111.

\bibitem{ref50} Holm Å, Tedgren Å, Larsson T. Heuristics for integrated optimization of catheter positioning and dwell time distribution in prostate HDR brachytherapy. ANNALS OF OPERATIONS RESEARCH. 2016;236(2):319-339.

\bibitem{ref51} Sachpazidis I, Hense J, Mavroidis P, Gainey M, Baltas D. Investigating the role of constrained CVT and CVT in HIPO inverse planning for HDR brachytherapy of prostate cancer. MEDICAL PHYSICS. 2019;46(7):2955-2968.

\bibitem{ref52} Gerlach S, Siebert FA, Schlaefer A. Robust stochastic optimization of needle configurations for robotic HDR prostate brachytherapy. MEDICAL PHYSICS. 2024;51(1):464-475.

\bibitem{ref53} Schulz JB, Bednarz BP, Floberg JM, Slagowski JM. Robust optimization of catheter positions and dwell times for HDR prostate brachytherapy. MEDICAL PHYSICS. 2026;53(9).

\bibitem{ref54} Zhang DG, Yang ZY, Jiang S, Zhou L, Zhou ZY, Wang W. Individualized and inverse optimized needle configuration for combined intracavitary-interstitial brachytherapy in locally advanced cervical cancer. JOURNAL OF CANCER RESEARCH AND THERAPEUTICS. 2019;15(7):1589-1596.

\bibitem{ref55} van der Meer MC, Bosman PAN, Niatsetski Y, et al. Bi-objective optimization of catheter positions for high-dose-rate prostate brachytherapy. MEDICAL PHYSICS. 2020;47(12):6077-6086.

\bibitem{ref56} Wang TH, Beaudry J, Núñez DA, et al. Automatic treatment planning using reinforcement learning for high-dose-rate prostate brachytherapy. Phys Med Biol. 2026;71(16).

\bibitem{ref57} Lin SQ, Yi OY, Gu HK, et al. Toward fully automated pre-implantation planning for cervical cancer brachytherapy: a template-guided multi-criteria optimization framework for catheter placement and dose distribution. Phys Med Biol. 2026;71(6).

\bibitem{ref58} Lei Y, Wang TH, Fu YB, et al. Catheter position prediction using deep-learning-based multi-atlas registration for high-dose rate prostate brachytherapy. MEDICAL PHYSICS. 2021;48(11):7261-7270.

\bibitem{ref59} Patil S, Pan J, Abbeel P, Goldberg K. Planning Curvature and Torsion Constrained Ribbons in 3D With Application to Intracavitary Brachytherapy. IEEE TRANSACTIONS ON AUTOMATION SCIENCE AND ENGINEERING. 2015;12(4):1332-1345.

\bibitem{ref60} Laan RC, Nout RA, Dankelman J, van de Berg NJ. MRI-driven design of customised 3D printed gynaecological brachytherapy applicators with curved needle channels. 3D Printing in Medicine. 2019;5(1):8.

\bibitem{ref61} Mao Z, Zhao HF. 3D-printed vaginal cylindrical template with curved needle channels in MRI-guided vaginal cuff brachytherapy. J Contemp Brachytherapy. 2023;15(6):422-431.

\bibitem{ref62} Zhang J, Luo CL, Zhang WJ, Yan JF, Yu L, Qiu J. Individual curved-needle interstitial template created using three-dimensional printing for brachytherapy for distal parauterine tumor recurrence. MEDICAL PHYSICS. 2025;52(1):596-604.

\bibitem{ref63} Padasdao B, Imanaka R, Podder TK, Konh B. Curvilinear catheter implantation in HDR prostate brachytherapy: feasibility study. MEDICAL PHYSICS. 2024;51(9):6332-6347.

\bibitem{ref64} de Vries M, Christianen M, Luthart L, de Vries KC, Kolkman-Deurloo IKK, van den Dobbelsteen JJ. Dosimetric benefits and preclinical performance of steerable needles in HDR prostate brachytherapy. MEDICAL ENGINEERING \& PHYSICS. 2024;128.

\bibitem{ref65} Straathof R, Meijaard JP, van Vliet-Pérez SM, et al. Multibody dynamic modeling of the behavior of flexible instruments used in cervical cancer brachytherapy. MEDICAL PHYSICS. 2024;51(5):3698-3710.

\bibitem{ref66} Straathof R, van Vliet-Pérez SM, Kolkman-Deurloo IKK, et al. Automated planning of curved needle channels in 3D printed patient-tailored applicators for cervical cancer brachytherapy. Phys Med Biol. 2024;69(23).

\bibitem{ref67} Fröhlich G, Agoston P, Lövey J, Polgár C, Major T. The Effect of Needle Number on the Quality of High-dose-rate Prostate Brachytherapy Implants. PATHOLOGY \& ONCOLOGY RESEARCH. 2010;16(4):593-599.

\bibitem{ref68} Chatigny PY, Bélanger C, Poulin É, Beaulieu L. Catheters and dose optimization using a modified CVT algorithm and multi-criteria optimization in prostate HDR brachytherapy. MEDICAL PHYSICS. 2022;49(10):6575-6587.

\bibitem{ref69} Guthier CV, Aschenbrenner KP, Müller R, Polster L, Cormack RA, Hesser JW. Real-time inverse high-dose-rate brachytherapy planning with catheter optimization by compressed sensing-inspired optimization strategies. Phys Med Biol. 2016;61(16):5956-5972.

\bibitem{ref70} Wang C, Gonzalez Y, Shen CY, Hrycushko B, Jia X. Simultaneous needle catheter selection and dwell time optimization for preplanning of high-dose-rate brachytherapy of prostate cancer. Phys Med Biol. 2021;66(5).

\bibitem{ref71} Frank CH, Ramesh P, Lyu Q, et al. Analytical HDR prostate brachytherapy planning with automatic catheter and isotope selection. MEDICAL PHYSICS. 2023;50(10):6525-6534.

\bibitem{ref72} Beck A, Teboulle M. A Fast Iterative Shrinkage-Thresholding Algorithm for Linear Inverse Problems. SIAM JOURNAL ON IMAGING SCIENCES. 2009;2(1):183-202.

\bibitem{ref73} Bélanger C, Aubin S, Lavallée MC, Beaulieu L. Simultaneous catheter and multicriteria optimization for HDR cervical cancer brachytherapy with a complex intracavity/interstitial applicator. MEDICAL PHYSICS. 2024;51(3):2128-2143.

\bibitem{ref74} Guthier CV, Damato AL, Hesser JW, Viswanathan AN, Cormack RA. A fast inverse treatment planning strategy facilitating optimized catheter selection in image-guided high-dose-rate interstitial gynecologic brachytherapy. MEDICAL PHYSICS. 2017;44(12):6117-6127.

\bibitem{ref75} Pötter R, Tanderup K, Kirisits C, et al. The EMBRACE II study: The outcome and prospect of two decades of evolution within the GEC-ESTRO GYN working group and the EMBRACE studies. Clinical and Translational Radiation Oncology. 2018;9:48-60.

\bibitem{ref76} Liu DC, Nocedal J. On the limited memory BFGS method for large scale optimization. Mathematical Programming. 1989;45(1):503-528.

\bibitem{ref77} Byrd RH, Lu P, Nocedal J, Zhu C. A Limited Memory Algorithm for Bound Constrained Optimization. SIAM J Sci Comput. 1995;16(5):1190--1208.

\bibitem{ref78} Kirkpatrick S, Gelatt CD, Vecchi MP. Optimization by Simulated Annealing. Science. 1983;220(4598):671-680.

\bibitem{ref79} Rivard MJ, Coursey BM, DeWerd LA, et al. Update of AAPM Task Group No. 43 Report: A revised AAPM protocol for brachytherapy dose calculations. MEDICAL PHYSICS. 2004;31(3):633-674.

\end{thebibliography}

\endgroup

\end{document}
